%%%%%%%%%%%%%%%%%%%%%%%%%%%%%%%%%%%%%%%%%%%%%%%%%%%%
% FRONT MATTER
%%%%%%%%%%%%%%%%%%%%%%%%%%%%%%%%%%%%%%%%%%%%%%%%%%%%

\documentclass[preprint,authoryear,12pt]{elsarticle}

% --------------------------------------------------
% Packages
% --------------------------------------------------
\usepackage{amsmath,amssymb,amsfonts}
\usepackage{booktabs}
\usepackage{natbib}
\usepackage{graphicx}
\usepackage{hyperref}
\usepackage{setspace}
\usepackage{geometry}

\geometry{margin=1in}

% Hyperref settings
\hypersetup{
    colorlinks = true,
    citecolor  = blue,
    linkcolor  = blue,
    urlcolor   = blue
}

% --------------------------------------------------
% Journal Info
% --------------------------------------------------
\journal{Energy Policy / Climate Policy (Draft Working Paper)}

% --------------------------------------------------
% TITLE AND AUTHORS
% --------------------------------------------------

\begin{document}

\begin{frontmatter}

\title{
Penalty-Based Investment Allocation Models for CPAT:\\
A Transparent Alternative to Logit-Based Share Rules
}
\author{}
%\author[wb]{Florent McIsaac, [Carolyn?], [Stephane?], [Unnada?], ,[Hector?], Stephen Stretton\corref{cor1}}
%\ead{sstretton@worldbank.org} \cortext[cor1]{Corresponding author.}\address[wb]{Fiscal Policy and Growth Global Department, World Bank Group}





% --------------------------------------------------
% ABSTRACT
% --------------------------------------------------


\newpage
\begin{abstract}
Investment allocation across power-generation technologies shapes future electricity supply,
emissions, and fiscal outcomes. Existing modelling traditions—from least-cost optimisation
models such as TIMES to accounting systems like LEAP and the multinomial logit rule in
CPAT—offer useful but incomplete tools for climate policy analysis. Optimisation models are too complex
for rapid country work, while accounting frameworks do not endogenously determine investment
responses. CPAT’s logit approach is practical but forces positive shares for all technologies and
can complicate policy dialogue.

This paper develops two transparent alternatives that preserve CPAT’s simplicity while improving
interpretability. A velocity-penalised share rule adds a quadratic cost to concentration,
yielding a closed-form allocation that can give zero shares to high-cost technologies. A weak
acceleration penalty moderates abrupt year-to-year changes, producing realistic transition paths
without fixing the long-run mix.

We integrate these elements with simple system-level cost adjustments—flexibility benefits for
gas, integration costs for variable renewables, and physical constraints—and derive analytical
expressions for total effective cost, steady-state shares, and annual update rules. Numerical
examples demonstrate how the combined model supports long-term strategy work, macro-fiscal
analysis, and developing-country power sector transitions. The framework offers a transparent,
spreadsheet-ready alternative to both optimisation models and logit-based rules, well suited for
integration into CPAT.

The appendices are mostly concerned with how to also consider the impact of a closed form systems cost on updated  investment and retirement rule, integrating the velocity-acceleration rule.
\end{abstract}



% --------------------------------------------------
% KEYWORDS / JEL 
% --------------------------------------------------

\begin{keyword}
Power-sector modelling; CPAT; Investment allocation; Velocity penalty; Acceleration penalty; LCOE; Fiscal analysis; Energy transition; Renewable energy; Developing countries.
\end{keyword}

\end{frontmatter}


\newpage
%%%%%%%%%%%%%%%%%%%%%%%%%%%%%%%%%%%%%%%%%%%%%%%%%%%%
% 1. HIGHLIGHTS
%%%%%%%%%%%%%%%%%%%%%%%%%%%%%
\section*{Highlights}
\begin{itemize}
    \item Introduces velocity- and acceleration-penalised allocation models as alternatives to logit rules.
    \item Provides analytical steady-state shares and transition paths suitable for CPAT.
    \item Embeds system costs (flexibility, integration, physical constraints) in a unified framework.
    \item Produces spreadsheet-ready formulae for total effective LCOE and update rules.
    \item Offers calibration guidance and policy-relevant interpretations for fiscal and LTS work.
\end{itemize}
\clearpage
\tableofcontents
\thispagestyle{empty}
\clearpage

%%%%%%%%%%%%%%%%%%%%%%%%%%%%%%%%%%%%%%%%%%%%%%%%%%%%
% 0. INTRODUCTION
%%%%%%%%%%%%%%%%%%%%%%%%%%%%%%%%%%%%%%%%%%%%%%%%%%%%

\section{Introduction}

The allocation of investment across competing power-generation technologies is one of the most
consequential decisions in long-term energy planning. It determines not only the cost,
reliability, and environmental footprint of the electricity system, but also the macroeconomic and
fiscal conditions under which countries pursue development and climate goals. Over the past
several decades, a wide range of modelling paradigms has emerged to support these decisions,
each reflecting a distinct view of how technology choices are made and how transitions unfold.
Some are grounded in economic optimisation, others in engineering accounting, others in
behavioural choice theory, and still others in evolutionary competition among technologies.
Understanding their respective strengths and weaknesses is a prerequisite for designing
allocation rules suitable for policy-focused tools such as the Climate Policy Assessment Tool
(\textsc{CPAT}).

The most widely used approach in energy-system planning is least-cost optimisation. Frameworks such as TIMES, MARKAL, MESSAGE or OSeMOSYS represent the energy system as a set of interlinked technologies, resources and constraints, and compute investment and dispatch paths that minimise the discounted cost of supplying energy. In simplified form, these models choose capacity for each technology and hourly generation levels so as to minimise total fixed and variable costs, subject to three core constraints: (i) hourly generation must meet hourly demand; (ii) generation from each technology cannot exceed its installed capacity adjusted by its availability in each hour; and (iii) total firm capacity must satisfy a planning reserve margin to ensure reliability.

The dual variables associated with the energy-balance and reliability constraints naturally define system-level marginal costs. These system costs—reflecting curtailment, flexibility limits, or reliability burdens—can be incorporated into reduced-form planning models as an additive “system cost” term appended to each technology’s private levelised cost \citep{grubb2019, miller2023}. This reduced-form mapping motivates the treatment of effective costs in models such as \textsc{CPAT}.

A second major paradigm is represented by accounting frameworks such as \textsc{LEAP}
\citep{leap2020} and \textsc{Sisepuede} \citep{sisepuede2024}, which emphasise transparent,
user-defined scenario construction rather than endogenous optimisation or behavioural
response. These systems excel at incorporating sectoral detail, stakeholder guidance and
narrative consistency. However, because technology shares are exogenously specified, they
cannot simulate the response of investment behaviour to changing policy or cost signals. For a
tool such as \textsc{CPAT}, which must trace how carbon prices, taxes or subsidies propagate
through investment patterns, this absence of behavioural structure remains a central
limitation \citep{CPATdocs}.

A contrasting behavioural approach is found in the evolutionary dynamics of the FTT:Power
framework \citep{Mercure2012}. Here, technologies compete for market share according to a
pairwise substitution mechanism shaped by relative costs, learning effects and path dependence.
In stylised form, the governing equation is:
\begin{equation}
\frac{ds_i}{dt} = \sum_{j \neq i} s_i \, s_j \, F_{ij}(C_j - C_i),
\end{equation}
where $s_i$ and $s_j$ are the \textit{shares of total capacity} held by technologies $i$ and $j$, 
$\frac{ds_i}{dt}$ is the corresponding \textit{change in the capacity share} of technology $i$, 
and $C_i$ and $C_j$ denote their respective generalised costs. The pairwise choice probability $F_{ij}$ is represented by a logit preference function:$\{F_{ij}(C_j - C_i) = 
\frac{1}{1 + \exp\!\left( \frac{C_i - C_j}{\sigma} \right)}\}$ where $\sigma$ captures the dispersion of investor preferences. This formulation produces S-curve diffusion,
lock-in, and path dependence.

Against this landscape, \textsc{CPAT} adopts a discrete-choice formulation based on the
multinomial logit model \citep{CPATdocs}, in which investment shares respond smoothly to
differences in levelised cost:
\begin{equation}
s_i
    = \frac{\exp(-C_i/\beta)}{\sum_j \exp(-C_j/\beta)}.
\label{eq:CPATlogit}
\end{equation}
This mapping is computationally trivial, transparent to calibrate and captures the intended
directional response of investment to policy shocks. In operational use within \textsc{CPAT},
$C_i$ is not simply the private levelised cost: it incorporates (i) \emph{system integration costs}
associated with variable renewables (derived from storage and balancing needs), (ii) a calibration
term capturing historical coal–gas relative competitiveness. In addition a hard scaling constraints on
annual renewable capacity additions (which reflect empirical construction bottlenecks) is imposed. These
adjustments enable \textsc{CPAT} to mimic, in reduced form, several features that would arise
endogenously in a more detailed hourly optimisation.

Recent work in climate economics introduces yet another perspective. Bauer, Hallegatte and
McIsaac \citeyearpar{BauerHallegatteMcIsaac2025} show that institutional or political frictions can
be represented via quadratic penalties that modify marginal abatement allocation. In a
single-period version of their framework, an optimal allocation of abatement $a_i$ across
sectors solves:
\begin{equation}
\min_{\{a_i\}} \sum_i \left( C_i a_i + \tfrac{\lambda}{2} a_i^2 \right)
    \quad \text{s.t. } \sum_i a_i = A,
\label{eq:hallegatte}
\end{equation}
yielding closed-form expressions that balance cost and institutional feasibility. This motivation
underpins the velocity- and acceleration-penalised allocation rules proposed in this paper.

Taken together, these modelling traditions motivate the search for an allocation mechanism that
preserves the interpretability and computational lightness of the logit rule, while allowing
corner solutions, representing diversification pressures and encoding simple transition
frictions. This paper develops precisely such mechanisms. The velocity-penalised model
produces a closed-form investment allocation in which high-cost technologies may receive zero
share, while the acceleration-penalised model limits abrupt year-to-year reallocations. Both
models can be integrated directly into \textsc{CPAT}, complementing its existing system-cost
adjustments and renewable scale-up constraints.

The remainder of the paper introduces a standard least cost optimisation model before velocity or acceleration effects. It then situates the two adjustment-cost alternatives within the broader modelling landscape, derives their analytical properties, illustrates their behaviour through numerical examples, and
discusses their relevance for fiscal and policy analysis in developing-country contexts, where transparency, interpretability and minimal data requirements are essential.



%%%%%%%%%%%%%%%%%%%%%%%%%%%%%%%%%%%%%%%%%%%%%%%%%%%%
% END OF PART 0 (FULL PROSE)
%%%%%%%%%%%%%%%%%%%%%%%%%%%%%%%%%%%%%%%%%%%%%%%%%%%%


%%%%%%%%%%%%%%%%%%%%%%%%%%%%%%%%%%%%%%%%%%%%%%%%%%%%
% 2. MODEL 1: VELOCITY-PENALISED INVESTMENT ALLOCATION
%%%%%%%%%%%%%%%%%%%%%%%%%%%%%%%%%%%%%%%%%%%%%%%%%%%%
\newpage
\section{Detailed Power Optimisation Models and their Representation In Principle in Reduced Form Models such as CPAT}

Before introducing the penalty-based investment rules, it is useful to clarify how the underlying
technological cost structure should be represented in a reduced-form model such as CPAT. In a
realistic power system, the cost of deploying an additional unit of coal, gas, solar or wind is not
simply its private levelised cost of electricity (LCOE). Instead it reflects a combination of energy
production, reliability contributions and balancing requirements, all of which are jointly
determined by the optimal dispatch of generators across the 8,760 hours of the year. Thus, any
parsimonious modelling framework must begin with an understanding of how these
system-level interactions translate into an effective marginal cost for each technology.

Consider a simplified generation expansion and dispatch model in which the planner determines
hourly output $q_{i,h}$ and installed capacity $K_i$ for each technology $i$ to minimise total
system cost subject to hourly load balance, capacity constraints and reserve or reliability
requirements. Formally, the planner solves:
\begin{equation}
    \min_{\{K_i, q_{i,h}\}}
       \sum_i \left( F_i K_i + V_i \sum_{h=1}^{8760} q_{i,h} \right)
\end{equation}
subject to:
\begin{align}
    &\sum_i q_{i,h} = L_h \qquad \forall h, \\
    &q_{i,h} \le K_i \cdot \text{CF}_{i,h} \qquad \forall i,h, \\
    &\sum_i \text{ELCC}_i(K_1,\dots,K_n)\, K_i \ge \text{PRM} \cdot \text{PeakLoad}.
\end{align}
Here $F_i$ and $V_i$ are fixed and variable costs, $\text{CF}_{i,h}$ is the availability profile, and
$\text{ELCC}_i$ denotes the effective load-carrying capability of technology $i$. The last
constraint ensures that the combined capacity provides sufficient reliability given the peak-load
requirement (PRM). This problem encompasses the basic operational and reliability features of
modern power systems and captures interactions between variable renewables and firm
capacity.

In this framework, the optimal solution produces not only the least-cost mix of capacities
$\{K_i\}$ but also a set of shadow prices on the constraints. These shadow prices have natural
interpretations as \emph{system costs} and \emph{system benefits}. For example:
\begin{itemize}
    \item The dual value on the hourly energy balance constraints indicates the marginal energy
          value of an additional unit of generation in each hour.
    \item The dual value on the firm-capacity (ELCC) constraint reflects the marginal value of
          reliability, which depends strongly on a technology's load-carrying capability.
    \item The binding of ramping, curtailment or flexibility constraints yields a marginal value
          associated with flexibility, which typically favours gas and storage.
\end{itemize}

When the model is solved for multiple candidate technologies, these shadow prices can be
combined into a \emph{virtual marginal cost} or \emph{system-adjusted cost} for each technology.
For example, the effective marginal cost of a unit of solar capacity may be written as:
\begin{equation}
    C^{\text{virt}}_{\text{solar}}
    = C^{\text{LCOE}}_{\text{solar}}
      + \gamma_{\text{rel}} \left(1 - \text{ELCC}_{\text{solar}}\right)
      + \gamma_{\text{flex}} \cdot \text{FlexReq}_{\text{solar}}
      + \gamma_{\text{cur}} \cdot \text{Curtailment}_{\text{solar}},
\end{equation}
where each $\gamma$ denotes the relevant shadow price from the optimisation and the terms
represent, respectively, the capacity credit shortfall, flexibility burden and expected curtailment.
Gas and coal would obtain corresponding adjustments, although the signs would differ:
flexibility tends to be a benefit for gas, while firm capacity is a benefit for coal or nuclear.

A key point is that these virtual costs reflect the system-level consequences of adding marginal
capacity of each technology—not the private LCOE alone. For a reduced-form model such as
CPAT, this provides a natural way to translate complex 8,760-hour optimisation outcomes into a
scalar \emph{system cost adder}:
\begin{equation}
    C_{i}^{\text{sys}} = C_{i}^{\text{virt}} - C_{i}^{\text{LCOE}}.
\end{equation}
The total effective cost for technology $i$ is then
\begin{equation}
    C_{i}^{\text{eff}} = C_{i}^{\text{LCOE}} + C_{i}^{\text{sys}}.
\end{equation}

In CPAT this effective marginal cost should be interpreted as a cost associated with the
\emph{capacity share} of each technology rather than the investment share. Once these
system-adjusted costs are obtained (e.g.~from a stylised or precomputed optimisation), they
serve as the baseline cost input to the velocity and acceleration models introduced in the next
section. In effect, the detailed optimisation maps into a reduced-form cost vector
$\{C_{i,t}^{\text{eff}}\}$ which captures the essential operational and reliability interactions.

The investment dynamics can then be interpreted as follows. In the absence of additional
constraints, a rational investor would allocate all investment to the technology with the lowest
$C_{i}^{\text{eff}}$. 

Please see Appendix A for a mathematical specifcation of how an overall 8760 hour optimisation model can specify a reduced form 1 time period model before the introduction of velocity and acceleration.


Diversification pressures, adjustment frictions and political-economy
constraints require smoothing of this allocation response. The velocity penalty introduces a systemic cost to excessive concentration, while the acceleration penalty limits rapid year-to-year
changes. Together, these penalties allow CPAT to generate both a plausible long-run allocation
based on effective system costs and realistic transition paths that respect operational and
institutional constraints. The next sections introduce 'clean' velocity and acceleration models and then introduce the combined system.


\newpage
%%%%%%%%%%%%%%%%%%%%%%%%%%%%%%%%%%%%%%%%%%%%%%%%%%%%
% 2. MODEL 1: VELOCITY-PENALISED INVESTMENT ALLOCATION
%%%%%%%%%%%%%%%%%%%%%%%%%%%%%%%%%%%%%%%%%%%%%%%%%%%%

\section{A Velocity-Penalised Model of Investment Allocation}

The first alternative considered in this paper is a static model in which the allocation of new
investment across technologies is determined by a trade-off between relative cost competitiveness
and an explicit penalty for excessive concentration. The term ``velocity'' is used in analogy with
adjustment-cost models in macroeconomics: just as high adjustment speed incurs a penalty in
capital accumulation models, here a high share of investment allocated to any one technology
incurs a quadratic cost. This structure captures the intuition that technological, institutional, or
system-integration factors may discourage placing a disproportionate share of new capacity in a
single option, even if that option has the lowest levelised cost of electricity (LCOE). It also
produces closed-form solutions that are transparent, intuitive, and compatible with the
spreadsheet-based architecture of CPAT.

\subsection{Formulation}

Suppose that a country must allocate one unit of new investment across $N$ generation
technologies indexed by $i$. Let $s_i$ denote the share allocated to technology $i$, and let $C_i$
represent its levelised cost. In a purely cost-driven model without frictions, least-cost selection
would dictate that all investment flows to the lowest-cost option. In practice, however, countries
rarely adopt such extreme strategies, either because diversification reduces risk, because grid
integration concerns arise at high penetration rates, or because investment pipelines cannot be
scaled arbitrarily fast.

To capture these considerations, we posit that the policymaker chooses investment shares
$\{s_i\}$ to minimise an augmented cost function:
\begin{equation}
    \sum_{i=1}^{N} \left( C_i s_i + \frac{\lambda}{2} s_i^2 \right),
    \label{eq:velocity_objective}
\end{equation}
subject to the simplex constraint
\begin{equation}
    \sum_{i=1}^{N} s_i = 1, \qquad s_i \ge 0.
    \label{eq:simplex_constraint}
\end{equation}
The parameter $\lambda > 0$ determines the severity of the diversification pressure. A large
value of $\lambda$ makes concentration costly and leads to more uniform investment shares,
whereas a small value of $\lambda$ makes the penalty negligible and approximates a least-cost
allocation rule.

This formulation is closely related to the quadratic penalty structures used by
\cite{BauerHallegatteMcIsaac2025} in modelling political economy constraints on mitigation
effort. It can also be interpreted as a reduced-form representation of system integration costs or
institutional frictions that increase disproportionately as a single technology dominates the
investment mix.

\subsection{Derivation of the closed-form allocation rule}

To derive the optimal allocation, consider the Lagrangian
\begin{equation}
    \mathcal{L}(s_1,\ldots,s_N;\mu)
    = \sum_{i=1}^{N} \left( C_i s_i + \frac{\lambda}{2} s_i^2 \right)
    - \mu \left( \sum_{i=1}^{N} s_i - 1 \right),
\end{equation}
where $\mu$ is the multiplier associated with the constraint
(\ref{eq:simplex_constraint}). The first-order conditions for an interior solution are
\begin{equation}
    \frac{\partial \mathcal{L}}{\partial s_i} = C_i + \lambda s_i - \mu = 0,
\end{equation}
which implies
\begin{equation}
    s_i = \frac{\mu - C_i}{\lambda}.
    \label{eq:s_i_expression}
\end{equation}
Summing over $i$ and enforcing the share constraint gives
\begin{equation}
    \sum_{i=1}^{N} s_i 
    = \frac{1}{\lambda}\left( N \mu - \sum_{i=1}^{N} C_i \right) = 1,
\end{equation}
leading to
\begin{equation}
    \mu = \bar{C} + \frac{\lambda}{N}, 
\end{equation}
where $\bar{C} = \frac{1}{N} \sum_{i=1}^{N} C_i$ denotes the average cost across technologies.
Substituting this back into (\ref{eq:s_i_expression}) yields the closed-form solution
\begin{equation}
    s_i
    = \frac{1}{N} + \frac{\bar{C} - C_i}{\lambda}.
    \label{eq:closedform_velocity}
\end{equation}

Expression (\ref{eq:closedform_velocity}) reveals several desirable properties. First, the model
naturally assigns higher shares to technologies with lower LCOE, with the magnitude of the
response governed by $1/\lambda$. Second, if a technology has a sufficiently high cost relative to
others, the right-hand side becomes negative, causing the optimal share to be zero once the
non-negativity constraint is enforced. Thus the model endogenously produces corner solutions in
which dominated technologies receive no investment, addressing a key limitation of the logit rule.
Third, the rule collapses to uniform diversification $s_i = 1/N$ as $\lambda \rightarrow \infty$,
and to least-cost allocation (a corner solution) as $\lambda \rightarrow 0$.

\subsection{Boundary behaviour and zero-share conditions}

The closed-form expression
(\ref{eq:closedform_velocity}) represents the interior solution. However, if
\begin{equation}
    \frac{1}{N} + \frac{\bar{C} - C_i}{\lambda} < 0,
\end{equation}
then the non-negativity constraint binds and $s_i = 0$. For such technologies, the remaining
shares renormalise to satisfy the simplex constraint. Because the algebra is fully transparent,
analysts can explicitly determine the cost threshold above which a technology is excluded. Let
$S = \sum_{j \ne i} C_j$ denote the sum of costs for all other technologies. Then the condition for
technology $i$ to retain a positive share can be expressed as
\begin{equation}
    C_i < \frac{S + \lambda}{N - 1}.
\end{equation}
This threshold determines the maximum cost a technology may have while still contributing to the
investment mix.

\subsection{Interpretation in terms of systemic LCOE}

The model also has a compelling economic interpretation. Differentiating the objective function
(\ref{eq:velocity_objective}) with respect to $s_i$ yields an effective or ``systemic'' marginal cost
\begin{equation}
    C_i^{\text{systemic}} = C_i + \lambda s_i.
\end{equation}
Thus the quadratic penalty can be understood as a technology-specific uplift applied in
proportion to the investment share already allocated to that technology. In systems with high
integration costs, for example, the systemic cost interpretation may provide a more meaningful
representation of planners' behaviour than either discrete-choice probabilities or optimisation
duality. It also offers a natural way to incorporate diversification pressures without imposing
exogenous constraints on maximum allowable shares.

\subsection{Calibration of the penalty parameter}

Calibrating $\lambda$ is central to applying this model in practice. If $\lambda$ is too small, the
model degenerates into a least-cost allocation rule, which may produce unrealistic concentration
in a single technology. If $\lambda$ is too large, the model forces a nearly uniform distribution
that ignores economically meaningful cost differences. In CPAT-like contexts, where investment
shares are updated annually and parameter transparency is important, one approach is to
calibrate $\lambda$ to match observed historical diversification patterns. Another approach is to
choose $\lambda$ such that a specified cost differential—for example, a 20 USD/MWh advantage—
translates into a desired difference in investment shares. Because the formula
(\ref{eq:closedform_velocity}) is linear in costs, such calibration is straightforward.

A further advantage is that the linear structure allows policymakers to inspect directly how
results vary with $\lambda$. Sensitivity analysis can therefore be communicated in terms that are
comprehensible to policymakers, who often prefer explicit trade-offs between diversification
pressure and cost responsiveness rather than the probabilistic responses implied by logit models.

\subsection{Advantages relative to CPAT’s multinomial logit rule}

Compared with the logit rule used in CPAT, the velocity-penalised model has several appealing
features. It produces zero shares for highly uncompetitive technologies, avoids the obligatory
interior solution of logit, offers direct economic interpretation via systemic LCOE, and provides
a cost–share relationship that is both linear and transparent. Whereas CPAT’s logit rule ensures
smoothness through an exponential function whose curvature may be unintuitive, the
velocity model achieves smoothness by adjusting the effective costs in a way that is easy to
explain to finance ministries.

%%%%%%%%%%%%%%%%%%%%%%%%%%%%%%%%%%%%%%%%%%%%%%%%%%%%
% END OF PART 2
%%%%%%%%%%%%%%%%%%%%%%%%%%%%%%%%%%%%%%%%%%%%%%%%%%%%
\newpage
%%%%%%%%%%%%%%%%%%%%%%%%%%%%%%%%%%%%%%%%%%%%%%%%%%%%
% 3. MODEL 2: ACCELERATION-PENALISED DYNAMICS
%%%%%%%%%%%%%%%%%%%%%%%%%%%%%%%%%%%%%%%%%%%%%%%%%%%%

\section{An Acceleration-Penalised Model of Transition Dynamics}

The second alternative developed in this paper recognises that investment decisions are rarely
made in isolation for a single year. Instead, they are embedded in multi-year planning cycles in
which existing project pipelines, policy commitments, and administrative capacities constrain the
pace at which a country can reconfigure its power-generation mix. For example, even if solar
emerges as the lowest-cost technology, the availability of transmission capacity, the structure of
public procurement, and the maturity of supply chains may limit the speed at which solar
investment can be scaled up. Conversely, politically sensitive sectors such as coal or gas may not
be reduced arbitrarily quickly because of employment concerns or regional development priorities. To capture these transition frictions, we introduce a model in which the investment decision is not penalised
for the \emph{level} of investment shares but for the \emph{change} in shares over time. Whereas the
velocity-penalised model produces an instantaneous steady-state allocation, the acceleration-
penalised model governs the path taken to move from one allocation to another.

\subsection{A dynamic investment allocation problem}

Let $s_{i,t}$ denote the share of investment in technology $i$ during year $t$. Let $C_{i,t}$ denote
the cost of that technology in year $t$, which may evolve due to learning, fuel prices, or policy.
We assume that the planner minimises the one-step cost
\begin{equation}
    C_{i,t} s_{i,t} + \frac{\lambda}{2} (s_{i,t} - s_{i,t-1})^2.
    \label{eq:acceleration_objective}
\end{equation}
The first term captures contemporaneous cost competitiveness, while the second term imposes a
quadratic penalty on deviations from last year's allocation. The penalty parameter $\lambda>0$
represents transition friction: large values of $\lambda$ imply that the system can adjust only
slowly, while small values allow rapid reallocation in response to cost signals.

Formally, at each time step the planner solves:
\begin{equation}
    \min_{\{s_{i,t}\}} 
    \sum_i \left( C_{i,t} s_{i,t} + \frac{\lambda}{2} (s_{i,t} - s_{i,t-1})^2 \right),
    \qquad \text{s.t.} \quad \sum_i s_{i,t} = 1,\; s_{i,t} \ge 0.
\end{equation}

This problem resembles dynamic adjustment-cost models in investment theory, where the cost of
changing the capital stock introduces inertia. Here, inertia applies to investment shares rather
than capacities, making the formulation suitable for spreadsheet environments and for annual
policy simulations such as those performed in CPAT.

\subsection{Derivation of the update rule}

To derive the first-order condition, consider the Lagrangian
\begin{equation}
    \mathcal{L}_t = \sum_i \left( C_{i,t} s_{i,t}
    + \frac{\lambda}{2} (s_{i,t} - s_{i,t-1})^2 \right)
    - \mu_t \left( \sum_i s_{i,t} - 1 \right),
\end{equation}
where $\mu_t$ enforces the share constraint in period $t$.

Differentiating with respect to $s_{i,t}$ yields:
\begin{equation}
    C_{i,t} + \lambda (s_{i,t} - s_{i,t-1}) - \mu_t = 0,
    \label{eq:acceleration_foc}
\end{equation}
so that
\begin{equation}
    s_{i,t} = s_{i,t-1} + \frac{\mu_t - C_{i,t}}{\lambda}.
    \label{eq:acceleration_update}
\end{equation}

Equation (\ref{eq:acceleration_update}) provides an elegant and transparent update rule. The
share of technology $i$ adjusts upward if its cost is below the shadow price $\mu_t$ and downward
if its cost is above it. The factor $1/\lambda$ determines the sensitivity of adjustment.

The multiplier $\mu_t$ ensures the shares sum to one. Summing
(\ref{eq:acceleration_update}) over all $i$ gives:
\begin{equation}
    1 = 1 + \frac{N \mu_t - \sum_i C_{i,t}}{\lambda},
\end{equation}
implying
\begin{equation}
    \mu_t = \frac{1}{N} \sum_i C_{i,t}.
\end{equation}
Thus, as in the velocity model, the shadow price equals the average cost across technologies.

Substituting into (\ref{eq:acceleration_update}) yields the final update rule:
\begin{equation}
    s_{i,t}
    = s_{i,t-1} 
    + \frac{\bar{C}_t - C_{i,t}}{\lambda}.
    \label{eq:s_final_acceleration}
\end{equation}

This rule is simple to implement and interpret: technologies that become cheaper than the
current average cost gain investment share, and those that become more expensive lose share,
with the speed of adjustment governed by $\lambda$.

\subsection{Interpretation via Euler–Lagrange dynamics}

The update rule (\ref{eq:s_final_acceleration}) admits an interpretation analogous to Euler–
Lagrange equations in continuous-time optimisation. If time were continuous and costs varied
smoothly, one would obtain a differential equation of the form:
\begin{equation}
    \lambda \frac{d s_i}{dt} = \bar{C}(t) - C_i(t).
\end{equation}
In this interpretation, the cost gradient acts as a ``force'' pushing the system toward more
cost-effective technologies, while the parameter $\lambda$ acts as a frictional term that 
dampens the response. Although our interest is in discrete annual steps, the analogy clarifies the
interpretation: the acceleration model moderates the velocity of change in the investment mix.

\subsection{Steady-state properties}

Unlike the velocity model, the acceleration model does not yield an interior steady-state share
distribution unless all technologies have identical costs. From
(\ref{eq:s_final_acceleration}), a steady state $s_{i,t} = s_{i,t-1}$ would require
\begin{equation}
    C_{i,t} = \bar{C}_t
    \quad \text{for all } i,
\end{equation}
which occurs only in degenerate cases. In practice, the model smoothly transitions toward a
corner solution determined by cost differences. This behaviour closely resembles dynamic
versions of the multinomial logit rule, where shares evolve gradually even if the long-run 
allocation collapses to the least-cost option.

In many policy contexts, the transition path is more relevant than the asymptotic distribution.
Governments often commit to long-term technology targets but are constrained by the pace at
which investment can realistically be shifted. The acceleration model provides a natural means of
representing such constraints.

\subsection{Handling boundary constraints}

If equation (\ref{eq:s_final_acceleration}) yields a negative value for $s_{i,t}$, the share is set to
zero and the remaining shares are renormalised. The model therefore permits both gradual
withdrawal from high-cost technologies and re-entry should costs fall in future years. This
feature is particularly useful for modelling gas, coal, or emerging technologies such as green
hydrogen, whose competitiveness can vary substantially over time.

\subsection{Relation to evolutionary diffusion models}

The structure of (\ref{eq:s_final_acceleration}) also bears a conceptual similarity to the
evolutionary diffusion equation underlying FTT:Power \cite{Mercure2012}. Both frameworks
govern the rate of change of market shares rather than the instantaneous level, and both embed a
cost gradient in the direction of substitution. However, the acceleration model is 
substantially simpler: it is linear rather than logistic, does not require pairwise substitution
probabilities, and produces easily interpretable parameters. Whereas FTT requires substantial
parameterisation and calibration, the acceleration model can be implemented in Excel with a
single free parameter.

\subsection{Calibration of transition speed}

The penalty parameter $\lambda$ determines the responsiveness of the investment mix. A large
value of $\lambda$ implies gradual adjustment, reflecting institutional, technical, or financial
constraints on reallocation. A small value implies rapid shifts, suitable for scenarios with strong
policy support, abundant administrative capacity, or highly modular technologies such as solar.

Calibration may be based on historical rates of change in technology investment shares,
on plausible annual build rates for relevant technologies, or on political considerations such as
the feasibility of phasing down coal over a prescribed period. Because the model is linear and the
parameter has direct mechanical interpretation, calibration can be communicated clearly to
policymakers.

\subsection{Advantages and limitations}

The acceleration-penalised model provides a parsimonious and intuitive method for shaping
transition pathways. It is particularly well suited for LTS, NDC updates, or CPAT analyses where
a policymaker wants to understand how quickly the investment mix may be able to shift under
different assumptions about institutional constraints. Its principal limitation is that it does not
provide a meaningful long-run equilibrium unless costs converge. For equilibrium analysis, the
velocity model of Section 2 is therefore preferred. In practice, both models can be used
jointly: the velocity model to determine a cost-consistent long-run allocation and the acceleration
model to govern the pace of movement toward it.

%%%%%%%%%%%%%%%%%%%%%%%%%%%%%%%%%%%%%%%%%%%%%%%%%%%%
% END OF PART 3
%%%%%%%%%%%%%%%%%%%%%%%%%%%%%%%%%%%%%%%%%%%%%%%%%%%%
\newpage
%%%%%%%%%%%%%%%%%%%%%%%%%%%%%%%%%%%%%%%%%%%%%%%%%%%%
% 4. NUMERICAL ILLUSTRATIONS
%%%%%%%%%%%%%%%%%%%%%%%%%%%%%%%%%%%%%%%%%%%%%%%%%%%%

\section{Numerical Illustrations}

To illustrate the behaviour of the two penalty-based allocation models, this section presents a
stylised numerical example involving five representative power-generation technologies. The
example is not intended to reproduce any particular national context; rather, its purpose is to
clarify the mechanics of the allocation rules and show how they differ from logit-based and
optimisation-based approaches. All numerical values are chosen for pedagogical clarity.

\subsection{Technology set and baseline costs}

Consider the following set of technologies and levelised costs, expressed in USD/MWh:
\begin{center}
\begin{tabular}{l c}
\toprule
Technology & $C_i$ (USD/MWh) \\
\midrule
Solar     & 20 \\
Gas       & 30 \\
Coal      & 40 \\
Wind      & 50 \\
Marine    & 100 \\
\bottomrule
\end{tabular}
\end{center}
The ordering reflects a common pattern in many emerging markets: solar is highly competitive,
gas is moderately priced, coal and wind are somewhat more expensive, and marine technologies
remain prohibitively costly for near-term deployment.

Throughout this section, we consider two values of the penalty parameter $\lambda$:
\begin{itemize}
    \item $\lambda = 100$, representing moderate diversification pressure,
    \item $\lambda = 200$, representing stronger diversification pressure.
\end{itemize}
These values are chosen for illustration and do not reflect a specific empirical calibration.

\subsection{Steady-state allocation under the velocity model}

Using the closed-form velocity allocation rule derived in Section 2,
\begin{equation}
    s_i = \frac{1}{N} + \frac{\bar{C} - C_i}{\lambda},
    \label{eq:velocity_rule_repeat}
\end{equation}
we obtain steady-state shares for the five technologies.

The average cost is:
\[
    \bar{C} = \frac{20 + 30 + 40 + 50 + 100}{5} = 48.
\]

Substituting values into (\ref{eq:velocity_rule_repeat}) with $\lambda = 100$ yields:
\[
    s_{\text{solar}} = 0.20 + \frac{48 - 20}{100} = 0.48,
\]
\[
    s_{\text{gas}}      = 0.20 + \frac{48 - 30}{100} = 0.38,
\]
\[
    s_{\text{coal}}     = 0.20 + \frac{48 - 40}{100} = 0.28,
\]
\[
    s_{\text{wind}}     = 0.20 + \frac{48 - 50}{100} = 0.18,
\]
\[
    s_{\text{marine}}   = 0.20 + \frac{48 - 100}{100} = -0.32.
\]

Because negative shares violate the non-negativity constraint, the marine share is set to zero and
the remaining shares are renormalised to sum to one. The resulting investment mix is:
\begin{center}
\begin{tabular}{l c}
\toprule
Technology & Share (\%) \\
\midrule
Solar       & 36.4 \\
Gas         & 28.8 \\
Coal        & 21.2 \\
Wind        & 13.6 \\
Marine      & 0.0 \\
\bottomrule
\end{tabular}
\end{center}

Two observations follow immediately. First, the velocity model naturally eliminates marine
generation, reflecting its extremely high cost. Second, despite solar being the least-cost option,
it does not absorb all investment; instead, diversification arises endogenously through the penalty
term. This stands in contrast to least-cost optimisation models, which would place all investment
in solar absent additional constraints, and to logit models, which would force marine to retain a
positive share even when its cost is uncompetitive.

\subsection{Sensitivity to the diversification penalty}

If we increase the penalty to $\lambda = 200$, cost differentials exert weaker influence. The
closed-form expression becomes:
\[
    s_i = 0.20 + \frac{\bar{C} - C_i}{200}.
\]

Substituting values yields:
\[
    s_{\text{solar}} = 0.20 + \frac{28}{200} = 0.34,
\]
\[
    s_{\text{gas}} = 0.20 + \frac{18}{200} = 0.29,
\]
\[
    s_{\text{coal}} = 0.20 + \frac{8}{200} = 0.24,
\]
\[
    s_{\text{wind}} = 0.20 - \frac{2}{200} = 0.19,
\]
\[
    s_{\text{marine}} = 0.20 - \frac{52}{200} = -0.06.
\]

After clipping and renormalisation, marine is again excluded. The remaining technologies show
greater uniformity than in the $\lambda = 100$ case. This illustrates a general property: as $\lambda$
increases, the system becomes less sensitive to cost differences and more evenly diversified.

\subsection{Cost threshold for participation}

The velocity model provides an analytic condition for determining whether a technology receives a
positive share. A technology $i$ retains a positive share if:
\begin{equation}
    C_i < \frac{S + \lambda}{N-1},
    \label{eq:threshold_condition_repeat}
\end{equation}
where $S$ is the sum of costs of all other technologies.

For marine generation,
\[
    S = 20 + 30 + 40 + 50 = 140.
\]
The threshold for participation is therefore:
\[
    C_{\text{marine}}^{\max} = \frac{140 + 100}{4} = 60,
\]
when $\lambda = 100$. Because the actual cost of marine is 100 USD/MWh, the model assigns it a
zero share. This is consistent with the intuition that no amount of diversification pressure can
justify investment in a technology that is so far outside the competitive range.

If $\lambda$ were increased to 300, the threshold would rise to:
\[
    C_{\text{marine}}^{\max} = \frac{140 + 300}{4} = 110,
\]
making marine marginally admissible. This illustrates how the diversification parameter shapes
the model’s willingness to tolerate expensive technologies.

\subsection{Transition dynamics under the acceleration model}

To illustrate how the acceleration-penalised model governs transitions, consider again the
cost vector $(20, 30, 40, 50, 100)$. Suppose the system begins with an equal allocation:
\[
    s_{i,0} = 0.20 \quad \text{for all } i.
\]
Using the update rule derived in Section 3,
\[
    s_{i,t} = s_{i,t-1} + \frac{\bar{C} - C_i}{\lambda},
\]
and setting $\lambda = 100$, the first-step adjustments are:
\[
    s_{\text{solar},1} = 0.20 + \frac{28}{100} = 0.48,
\]
\[
    s_{\text{gas},1} = 0.20 + \frac{18}{100} = 0.38,
\]
\[
    s_{\text{coal},1} = 0.20 + \frac{8}{100} = 0.28,
\]
\[
    s_{\text{wind},1} = 0.20 - \frac{2}{100} = 0.18,
\]
\[
    s_{\text{marine},1} = 0.20 - \frac{52}{100} = -0.32.
\]

After clipping, the transition mirrors exactly the velocity allocation in the first period. In
subsequent years, the system continues to shift share toward solar and gas until marine's share
stabilises at zero. For values of $\lambda$ larger than 100, the transition occurs more slowly and
with more gradual year-to-year adjustments.

An important distinction emerges: under the velocity model, the above distribution is the
steady-state allocation; under the acceleration model, the system continues to evolve beyond the
first period unless cost differentials vanish. Thus the acceleration model is better interpreted as a
pathway generator rather than a steady-state predictor.

\subsection{Policy interpretation}

These numerical examples underline several implications for policy modelling. First, the velocity
model provides an analytically clean and intuitive formula for the distribution of investment in a
hypothetical steady state given relative costs and diversification pressure. Second, the acceleration
model provides a mechanism for exploring realistic short- and medium-run pathways, especially in
contexts where political or administrative constraints limit the pace of change. Third, the analytic
threshold for zero shares is particularly valuable in fiscal policy contexts, where the objective is
often to illustrate how taxes or subsidies can render technologies marginal or obsolete.

Both models are sufficiently transparent that their behaviour can be explained directly to finance
ministries, which considerably enhances their usefulness in tools such as CPAT.

%%%%%%%%%%%%%%%%%%%%%%%%%%%%%%%%%%%%%%%%%%%%%%%%%%%%
% END OF PART 4
%%%%%%%%%%%%%%%%%%%%%%%%%%%%%%%%%%%%%%%%%%%%%%%%%%%%
\newpage
\section{Land-use–related system costs and full effective cost}

In addition to integration and flexibility effects, large-scale deployment of variable renewables
is constrained by available land. We introduce a simple land-use–related system cost that grows
quadratically with the total land fraction occupied by solar and wind, and is already noticeable
when renewables occupy about $2.5\%$ of the country’s land area.

Let $A_{\text{tot}}$ denote total land area of the country (m$^2$). For solar PV, assume a peak
power density of $100$ W$_{\text{peak}}$/m$^2$ and a capacity factor of $20\%$. For wind, assume
$16$ W$_{\text{peak}}$/m$^2$ and a capacity factor of $25\%$. The capacity–area relationships
(in MW) are then:
\begin{align*}
    K_{\text{sol},t} &= \frac{100 \; A_{\text{sol},t}}{10^6}
                     = \frac{A_{\text{sol},t}}{10^4},
    \qquad \Rightarrow \quad
    A_{\text{sol},t} = 10^4\,K_{\text{sol},t},\\[4pt]
    K_{\text{win},t} &= \frac{16 \; A_{\text{win},t}}{10^6}
                     = \frac{A_{\text{win},t}}{6.25\times 10^4},
    \qquad \Rightarrow \quad
    A_{\text{win},t} = 6.25\times 10^4\,K_{\text{win},t},
\end{align*}
where $A_{\text{sol},t}$ and $A_{\text{win},t}$ are the land areas used by solar and wind,
respectively.

The total land area used by these renewables in year $t$ is
\[
    A_{\text{RES},t}
    = A_{\text{sol},t} + A_{\text{win},t}
    = 10^4 K_{\text{sol},t} + 6.25\times 10^4 K_{\text{win},t},
\]
and the corresponding land-use fraction is
\[
    f_t = \frac{A_{\text{RES},t}}{A_{\text{tot}}}.
\]

We assign the land-use penalty primarily to solar, with the cost scaling quadratically in the
total land fraction $f_t$. To make the effect noticeable, we choose the coefficient so that the
land-use cost for solar reaches approximately \$10/MWh when renewables occupy $2.5\%$ of
total land area:
\[
    C_{\text{sol},t}^{\text{land}}
    = \zeta_{\text{sol}} f_t^2,
    \qquad
    \zeta_{\text{sol}}
    = \frac{10}{0.025^2}
    \approx 1.6\times 10^4 \;\text{USD/MWh}.
\]
Thus, at $f_t = 0.025$,
\[
    C_{\text{sol},t}^{\text{land}} \approx 10 \;\text{USD/MWh},
\]
and the penalty grows rapidly as the land fraction increases. For wind, an analogous but
smaller coefficient $\zeta_{\text{win}}$ can be used if desired; for simplicity, it may be set to zero
or a fraction of $\zeta_{\text{sol}}$.

\medskip

Collecting all components, the full effective cost per MWh that drives the investment rule in
year $t$ can be written as
\begin{equation}
    C_{i,t}^{\text{eff,full}}
    = C_{i,t}^{\text{base}}
      + C_{i,t}^{\text{sys}}(s_{t-1};D)
      + C_{i,t}^{\text{land}}(f_t)
      + \lambda_v s_{i,t}
      + \lambda_a (s_{i,t} - s_{i,t-1}),
    \label{eq:full_effective_cost_land}
\end{equation}
where:
\begin{itemize}
    \item $C_{i,t}^{\text{base}}$ is the base LCOE;
    \item $C_{i,t}^{\text{sys}}(s_{t-1};D)$ captures integration, flexibility and other system
          effects (e.g.\ gas flexibility benefits, renewable balancing costs, density-related
          constraints);
    \item $C_{i,t}^{\text{land}}(f_t)$ is the land-use penalty defined above, significant when
          the combined solar–wind land fraction $f_t$ reaches a few percent of total area;
    \item $\lambda_v s_{i,t}$ and $\lambda_a (s_{i,t} - s_{i,t-1})$ are the marginal contributions
          of the velocity and acceleration penalties, respectively, arising from the quadratic
          investment objective.
\end{itemize}

The velocity–acceleration optimality conditions are obtained as in the main text, replacing
$C_{i,t}^{\mathrm{eff}}$ with $C_{i,t}^{\text{eff,full}}$ in the first-order conditions and in the
closed-form investment share rule.

\medskip

%%%%%%%%%%%%%%%%%%%%%%%%%%%%%%%%%%%%%%%%%%%%%%%%%%%%
% 6. CPAT IMPLEMENTATION RECOMMENDATIONS AND COMBINED SPECIFICATION
%%%%%%%%%%%%%%%%%%%%%%%%%%%%%%%%%%%%%%%%%%%%%%%%%%%%
\newpage
\section{CPAT Implementation Recommendations and Combined Specification}

This section offers concrete recommendations for implementing the proposed allocation models in
CPAT and presents a combined specification that incorporates both velocity and acceleration
penalties alongside simple representations of system benefits and costs. The aim is to retain
spreadsheet-level simplicity while capturing the most important features of realistic power system
transitions.

\subsection{Recommended structure for CPAT}

For CPAT and similar fiscal tools, it is natural to use the velocity-penalised model as the
backbone of the investment allocation module. The closed-form steady-state solution is
transparent, easy to explain to finance ministries, and provides a direct link between relative
costs, diversification pressure and the exclusion of high-cost technologies. This makes it well
suited for the primary counterfactual comparisons that CPAT performs across tax and subsidy
scenarios.

At the same time, it is often useful to moderate the speed at which investment shares respond to
policy-induced cost changes. For example, a sudden increase in carbon taxation might, in
principle, justify an immediate shift away from coal, but in practice administrative, technical and
political constraints will slow the response. For this reason, CPAT can include a weak
acceleration penalty that constrains year-to-year changes in investment shares without
dominating the allocation outcome. In practical terms, this means choosing the velocity penalty
parameter to govern the long-run distribution of investment and setting the acceleration penalty
parameter to a relatively small value that primarily affects short-run dynamics.


\subsection{A combined velocity–acceleration model}

Let $i$ index technologies and $t$ index years. Denote by $s_{i,t}$ the share of new investment in
technology $i$ in year $t$. For CPAT implementation we propose the following generic
one-period problem:
\begin{equation}
    \min_{\{s_{i,t}\}}
    \sum_i \left[
        C_{i,t}^{\text{base}} s_{i,t}
      + C_{i,t}^{\text{sys}}(s_{t-1};D) s_{i,t}
      + \frac{\lambda_v}{2} s_{i,t}^2
      + \frac{\lambda_a}{2}(s_{i,t} - s_{i,t-1})^2
    \right],
    \label{eq:combined_objective}
\end{equation}
subject to
\begin{equation}
    \sum_i s_{i,t} = 1, \qquad s_{i,t} \ge 0.
\end{equation}

Here, $C_{i,t}^{\text{base}}$ is the base levelised cost of electricity for technology $i$ in year $t$
(as already computed in CPAT), while $C_{i,t}^{\text{sys}}(s_{t-1};D)$ represents a set of
system-level adjustments evaluated on the previous year's shares $s_{t-1}$ and a country-level
measure of population density $D$. Evaluating system adjustments on lagged shares keeps the
problem linear-quadratic in $s_{t}$ and preserves analytical tractability.

In line with CPAT’s focus on a small number of interpretable terms, we propose the following
decomposition of $C_{i,t}^{\text{sys}}(s_{t-1};D)$:
\begin{equation}
    C_{i,t}^{\text{sys}}(s_{t-1};D)
    = - \beta_G G_i s_{\text{gas},t-1}
      + \kappa_R R_i S_{\text{RES},t-1}
      + \frac{\phi_i}{D} s_{i,t-1},
    \label{eq:C_sys_definition}
\end{equation}
where:
\begin{itemize}
    \item $G_i$ is an indicator equal to one if technology $i$ is gas, and zero otherwise; the term
          $-\beta_G G_i s_{\text{gas},t-1}$ captures the system \emph{benefit} of flexible gas
          capacity, which increases as the gas share rises, effectively lowering its marginal cost;
    \item $R_i$ is an indicator for variable renewable technologies (such as solar and wind); the
          term $\kappa_R R_i S_{\text{RES},t-1}$ captures additional system \emph{costs} associated
          with storage, balancing and grid integration that rise with the aggregate share of
          variable renewables
          \[
              S_{\text{RES},t-1} = \sum_j R_j s_{j,t-1};
          \]
    \item $\phi_i$ is a technology-specific parameter scaling a quadratic physical constraint
          linked to population density; the term $(\phi_i / D) s_{i,t-1}$ represents the idea that,
          for a given technology, the effective marginal system cost rises with its local spatial
          footprint when population density is low (for example, for land-intensive renewables in
          sparsely populated countries).
\end{itemize}

Substituting (\ref{eq:C_sys_definition}) into (\ref{eq:combined_objective}) and noting that the
system adjustments depend only on lagged shares, the objective remains quadratic in $s_{i,t}$. The
Lagrangian for year $t$ can therefore be written as
\begin{align}
    \mathcal{L}_t
    &= \sum_i \left[
        \left( C_{i,t}^{\text{base}} + C_{i,t}^{\text{sys}}(s_{t-1};D) \right) s_{i,t}
      + \frac{\lambda_v}{2} s_{i,t}^2
      + \frac{\lambda_a}{2}(s_{i,t} - s_{i,t-1})^2
    \right]
    \nonumber\\
    &\quad - \mu_t \left( \sum_i s_{i,t} - 1 \right).
\end{align}

The first-order condition with respect to $s_{i,t}$ is
\begin{equation}
    C_{i,t}^{\text{base}} + C_{i,t}^{\text{sys}}(s_{t-1};D)
    + \lambda_v s_{i,t}
    + \lambda_a (s_{i,t} - s_{i,t-1})
    - \mu_t = 0.
\end{equation}
Collecting terms in $s_{i,t}$ yields:
\begin{equation}
    (\lambda_v + \lambda_a) s_{i,t}
    = \mu_t
      - C_{i,t}^{\text{base}}
      - C_{i,t}^{\text{sys}}(s_{t-1};D)
      + \lambda_a s_{i,t-1}.
    \label{eq:combined_foc}
\end{equation}
Thus the optimal share in year $t$ is given by
\begin{equation}
    s_{i,t}
    = \frac{1}{\lambda_v + \lambda_a}
      \left[
          \mu_t
        - C_{i,t}^{\text{base}}
        - C_{i,t}^{\text{sys}}(s_{t-1};D)
        + \lambda_a s_{i,t-1}
      \right].
    \label{eq:combined_update_pre_mu}
\end{equation}

The multiplier $\mu_t$ is determined by the share constraint. Summing
(\ref{eq:combined_update_pre_mu}) over $i$ and imposing $\sum_i s_{i,t} = 1$ gives:
\begin{equation}
    1
    = \frac{N \mu_t
        - \sum_i C_{i,t}^{\text{base}}
        - \sum_i C_{i,t}^{\text{sys}}(s_{t-1};D)
        + \lambda_a \sum_i s_{i,t-1}
      }{\lambda_v + \lambda_a}.
\end{equation}
Using $\sum_i s_{i,t-1} = 1$, this simplifies to
\begin{equation}
    \mu_t
    = \bar{C}_t^{\text{base}}
      + \bar{C}_t^{\text{sys}}(s_{t-1};D)
      - \frac{\lambda_a}{N}
      + \frac{\lambda_v + \lambda_a}{N},
\end{equation}
where
\[
    \bar{C}_t^{\text{base}} = \frac{1}{N} \sum_i C_{i,t}^{\text{base}}, 
    \qquad
    \bar{C}_t^{\text{sys}}(s_{t-1};D) = \frac{1}{N} \sum_i C_{i,t}^{\text{sys}}(s_{t-1};D).
\]

It is convenient to define an \emph{adjusted average cost}
\begin{equation}
    \bar{C}_t^{\text{adj}}
    = \bar{C}_t^{\text{base}} + \bar{C}_t^{\text{sys}}(s_{t-1};D) - \frac{\lambda_a}{N}.
\end{equation}
Then the multiplier becomes
\begin{equation}
    \mu_t = \bar{C}_t^{\text{adj}} + \frac{\lambda_v + \lambda_a}{N}.
\end{equation}

Substituting this back into (\ref{eq:combined_update_pre_mu}) yields a closed-form update rule:
\begin{equation}
    s_{i,t}
    = \frac{1}{N}
      + \frac{
            \bar{C}_t^{\text{adj}}
          - C_{i,t}^{\text{base}}
          - C_{i,t}^{\text{sys}}(s_{t-1};D)
          + \lambda_a s_{i,t-1}
        }{\lambda_v + \lambda_a}.
    \label{eq:combined_update_final}
\end{equation}

Equation (\ref{eq:combined_update_final}) is the central expression for CPAT implementation. It
states that each technology’s share is the sum of (i) an equal-share baseline $1/N$ and (ii) an
adjustment term that depends on the difference between adjusted average cost and its own
effective cost, plus a contribution from last year’s share scaled by the acceleration parameter.
When $\lambda_a$ is set to zero and the system cost adjustments are dropped, this equation
collapses to the pure velocity model. When $\lambda_v$ is large and $\lambda_a$ is small, the
model behaves like a lightly smoothed version of least-cost allocation with system adjustments.

The total effective cost of investing in technology $i$ in year $t$ can be written compactly as:
\begin{equation}
    C_{i,t}^{\text{tot}}(s_t,s_{t-1};D)
    = C_{i,t}^{\text{base}}
      + C_{i,t}^{\text{sys}}(s_{t-1};D)
      + \lambda_v s_{i,t}
      + \lambda_a (s_{i,t} - s_{i,t-1}),
    \label{eq:total_LCOE}
\end{equation}
which clarifies the sense in which velocity and acceleration penalties adjust the levelised cost to
reflect diversification pressures and transition frictions.

\subsection{Calibration Strategy}

A practical calibration strategy for CPAT can proceed in stages. First, the system adjustment
parameters $(\beta_G, \kappa_R, \phi_i)$ can be informed by engineering studies and existing
estimates of balancing costs, storage requirements and grid expansion needs. These estimates can
be translated into approximate marginal cost adjustments per unit of gas flexibility, per unit of
variable renewable share and per unit of technological footprint scaled by population density.

Second, the velocity and acceleration parameters $(\lambda_v, \lambda_a)$ can be calibrated to
historical data on investment shares and costs. Specifically, one can collect a panel of
$\{C^{\text{base}}_{i,t}, s_{i,t}\}$ for a country or peer group and estimate $\lambda_v$ and
$\lambda_a$ so that the model-implied shares given by (\ref{eq:combined_update_final})
approximate observed shares. A simple approach is to minimise the sum of squared deviations
between observed and modelled shares over a historical window, possibly using regularisation to
ensure that $\lambda_a$ remains small relative to $\lambda_v$.

Finally, calibration results should be cross-checked for plausibility in policy dialogue. If the
estimated $\lambda_v$ implies almost uniform diversification regardless of costs, or if $\lambda_a$
implies unrealistically slow transitions, parameters can be adjusted within ranges consistent with
historical experience and stakeholder judgement. Because the parameters have direct mechanical
interpretation, these adjustments can be communicated explicitly, which is an advantage relative
to more opaque functional forms.

\subsection{Integration into CPAT}

CPAT requires a simple, robust, and transparent method for determining investment responses to
policy-induced cost changes. The existing multinomial logit rule provides smoothness and ease of
calibration but also forces positive shares for all technologies. The velocity model can replace or
complement this rule in applications where it is important to illustrate the complete elimination
of high-cost technologies or to ensure that investment shares reflect plausible diversification
patterns.

The acceleration model plays a distinct role. In many country engagements, the central
question is not which technologies ultimately dominate, but how quickly the system can transition.
CPAT scenarios that incorporate carbon taxation, fossil-fuel subsidy reform, or renewable energy
support policies can benefit from modelling the pace of reallocation rather than relying solely on
instantaneous responses. The acceleration model provides this capability while preserving the
interpretability required for communication with finance ministries.

Both models can be integrated into CPAT in a modular fashion. The velocity model defines a
steady-state target allocation for each policy scenario. The acceleration model can then govern
the multi-year path toward that target. This combination mirrors common practice in long-term
strategy (LTS) planning, where a country commits to a 2050 mix and then explores feasible
transition pathways.
%%%%%%%%%%%%%%%%%%%%%%%%%%%%%%%%%%%%%%%%%%%%%%%%%%%%
% 5. DISCUSSION AND POLICY INTEGRATION
%%%%%%%%%%%%%%%%%%%%%%%%%%%%%%%%%%%%%%%%%%%%%%%%%%%%
\newpage

\section{Discussion}

The preceding sections developed two analytically tractable alternatives to the multinomial logit
investment rule widely used in policy simulation tools such as CPAT. The velocity-penalised
model produces a closed-form steady-state distribution of investment shares that depends on
relative costs and a diversification parameter. The acceleration-penalised model governs the pace
of change in the investment mix, yielding smooth transitions consistent with institutional and
political constraints. This section situates these approaches within the broader modelling
landscape and outlines their implications for fiscal policy and long-term planning.

\subsection{Comparison with existing modelling paradigms}

Least-cost optimisation models such as TIMES represent the most rigorous and detailed approach
to power-sector investment planning. Their strength lies in the ability to encode a wide range of
technical, network, and resource constraints. However, their computational requirements render
them unsuitable for country-level fiscal diagnostics, where policymakers require rapid scenario
assessment in a spreadsheet environment. Moreover, least-cost optimisation tends to produce
corner solutions unless diversification constraints are explicitly introduced. The velocity model
developed here embeds diversification directly in the objective function, achieving a similar
effect with minimal computational overhead.

LEAP-style accounting models offer transparency and accessibility, but they lack explicit
behavioural rules determining investment allocation. Their strength lies in narrative scenario
construction rather than equilibrium modelling. The penalty-based models presented here provide
a behavioural closure that remains consistent with LEAP-style narrative pathways and can be
embedded directly in tools designed for policy communication.

FTT diffusion models offer a rich representation of technological substitution driven by pairwise
comparisons and learning-by-doing. They capture momentum effects, lock-in, and non-linear
dynamics, but they require large parameter sets and careful calibration. The acceleration model
shares with FTT the focus on transition dynamics, but its linear structure and single-parameter
governance make it much more suitable for fiscal policy tools. In contrast to logit-based models,
the penalty models allow a technology to drop out entirely when it is prohibitively costly.



\subsection{Macro-fiscal implications and MFMod integration}

From a fiscal perspective, investment allocation affects public revenues and expenditures
through several channels: electricity tariffs, fossil-fuel tax bases, subsidy requirements,
public investment obligations, and external financing needs. MFMod, the World Bank’s standard
macro-fiscal model, links sectoral investment to macroeconomic aggregates, allowing analysts to
quantify how different technology pathways affect debt sustainability, GDP growth, and fiscal
space.

The penalty-based approaches described here provide a lightweight but disciplined behavioural
rule that can be fed into MFMod. For example, an acceleration-constrained scenario in which
renewable energy uptake proceeds slowly may imply a prolonged period of fossil-fuel tax revenue
or sustained fuel import bills. In contrast, a velocity-based allocation with strong diversification
pressure may shift investment more rapidly toward low-cost technologies, reducing the fiscal
burden of subsidies or grid expansion.

Because the penalty parameters are interpretable, they can be mapped to institutional realities:
permitting timelines, grid bottlenecks, procurement cycles, or political feasibility constraints.
This mapping enables transparent and policy-relevant macro-fiscal scenarios.

\subsection{Relevance for long-term strategies and NDC updates}

Long-term strategies (LTS) and enhanced NDCs require countries to articulate realistic and
credible transition pathways. In these contexts, the penalty models offer several advantages. The
velocity model clarifies the cost context in which diversification or technology neutrality may be
desirable. The acceleration model allows analysts to constrain the pace of coal phase-out or
renewable scale-up in ways consistent with administrative and financial capacity.

For developing countries, where rapid transitions may be limited by fiscal constraints, these
models allow a direct representation of the trade-off between ambition and feasibility. They also
provide a principled means to evaluate the impact of international support, such as concessional
finance or carbon contracts for difference (CCfDs), on the speed and direction of investment
flows.

\subsection{Limitations and avenues for extension}

The penalty-based models deliberately abstract from several important features of real-world
power systems. They do not incorporate learning curves, resource depletion, reliability
constraints, grid integration costs, or storage requirements. Nor do they consider feed-in
constraints or investment indivisibilities. These features can be introduced in future extensions by
modifying the cost term or by adding constraints to the objective function.

The linearity of the acceleration model is both a strength and a limitation. It enables clear
interpretation but does not capture the logistic behaviour observed in historical technology
diffusion. Hybrid models that combine linear penalties with logit-like diffusion functions may
offer a pragmatic compromise between realism and tractability.

\newpage
%%%%%%%%%%%%%%%%%%%%%%%%%%%%%%%%%%%%%%%%%%%%%%%%%%%%
% END OF SECTION 6
%%%%%%%%%%%%%%%%%%%%%%%%%%%%%%%%%%%%%%%%%%%%%%%%%%%%
\newpage
%%%%%%%%%%%%%%%%%%%%%%%%%%%%%%%%%%%%%%%%%%%%%%%%%%%%
% 7. CONCLUSION
%%%%%%%%%%%%%%%%%%%%%%%%%%%%%%%%%%%%%%%%%%%%%%%%%%%%

\section{Recommendations}

The paper has the following two main recommendations for CPAT.

First adopt the velocity rule above as the main capacity constraint, replacing the renewable energy constraint. Include a minor acceleration term as well, so the overall equation includes systems cost, physical constraints, velocity and acceleration. Use the velocity rule instead of existing multilogit plus RE constraint.

Second, consider a better approach to useing systems costs to determine least cost optimised levels of capacity additions. Consider updating the investment and retirement rules to better reflect what that optimised levels of capacity is. The appendix modifies the CPAT approach in more detail, integrated with the velocity-acceleration rule.


\section{Conclusion}

This paper developed two analytically tractable alternatives to the multinomial logit rule
commonly used in power-sector investment modelling. The velocity-penalised model yields a
closed-form steady-state distribution of investment shares driven by relative costs and a
diversification parameter. The acceleration-penalised model governs the pace of transition,
capturing institutional and political constraints.

Both models are simple enough to be implemented in spreadsheets while retaining sufficient
structure for integration into CPAT and macro-fiscal tools such as MFMod. Their transparency
makes them particularly suitable for engagement with finance ministries, where clarity and
interpretability are essential. By providing both steady-state and dynamic behaviour, the models
support long-term planning, fiscal risk analysis, and the design of realistic policy pathways.

Future work may extend these models to include learning dynamics, resource depletion,
explicit grid constraints, or hybrid functional forms blending diversification penalties with
diffusion-based approaches. Nonetheless, even in their present form, the models offer a valuable
addition to the analytical toolkit for climate and fiscal policy.
%%%%%%%%%%%%%%%%%%%%%%%%%%%%%%%%%%%%%%%%%%%%%%%%%%%%
% REFERENCES
%%%%%%%%%%%%%%%%%%%%%%%%%%%%%%%%%%%%%%%%%%%%%%%%%%%%

\bibliographystyle{plainnat}   % Elsevier author-year style
\bibliography{references}        % References stored 

%%%%%%%%%%%%%%%%%%%%%%%%%%%%%%%%%%%%%%%%%%%%%%%%%%%%
% ACKNOWLEDGEMENTS (Optional)
%%%%%%%%%%%%%%%%%%%%%%%%%%%%%%%%%%%%%%%%%%%%%%%%%%%%
\section{Acknowledgements}
To be added later. Paper written with substantial input from Florent McIsaac.

%%%%%%%%%%%%%%%%%%%%%%%%%%%%%%%%%%%%%%%%%%%%%%%%%%%%
% DECLARATIONS (Optional)
%%%%%%%%%%%%%%%%%%%%%%%%%%%%%%%%%%%%%%%%%%%%%%%%%%%%
\section*{Declaration of Interests}
The author declares no conflicts of interest.

\section*{Data Availability}
No new data were created for this study. All modelling examples use stylised inputs.

\newpage
%%%%%%%%%%%%%%%%%%%%%%%%%%%%%%%%%%%%%%%%%%%%%%%%%%%%
% APPENDIX A: MATHEMATICAL DERIVATIONS
%%%%%%%%%%%%%%%%%%%%%%%%%%%%%%%%%%%%%%%%%%%%%%%%%%%%
\appendix

\newpage
\section*{Introduction to Appendices}
The appendices provides derivations and further work omitted from the main text for readability. Whilst the focus of the main paper is on the velocity and acceleration models, here the focus is relating the capacity investment rule relative to optimised levels or, assuming those are not available, to the stylised systems costs. As such we first look at the reduced form of the optimisation model, then the investment rule with known (least) cost capacities known, and then finally the investment rule with only (marginal) systems costs known.

A final appendix discusses the algorithmic implementation in Excel.


\section{Yearly Reduced form of Least Cost Hourly Optimisation Models without Velocity or Acceleration Penalties}


\subsection{Static least-cost benchmark and shadow values.}

Consider a stylised least-cost capacity and dispatch problem over 8760 hours.
The planner chooses capacities $K_i$ and hourly outputs $q_{i,h}$ to minimise
annualised fixed and variable costs:
\begin{equation}
\begin{aligned}
\min_{\{K_i, q_{i,h}\}} \quad 
 & \sum_i F_i K_i \;+\; \sum_{h=1}^{8760} \sum_i V_i q_{i,h} \\
\text{s.t.} \quad 
 & \sum_i q_{i,h} = L_h, && \forall h, \\
 & 0 \le q_{i,h} \le K_i CF_{i,h}, && \forall i,h, \\
 & \sum_i \mathrm{ELCC}_i K_i \;\ge\; \mathrm{PRM}\cdot \mathrm{PeakLoad}. &&
\end{aligned}
\end{equation}
Let $\lambda_h$ denote the dual variables on the hourly balance constraints
(energy shadow prices), $\phi_{i,h}\ge 0$ the duals on the upper bounds
$q_{i,h} \le K_i CF_{i,h}$ (capacity scarcity in each hour), and
$\mu \ge 0$ the dual on the reliability constraint. The Lagrangian is
\begin{equation}
\begin{aligned}
\mathcal{L} \;=\;
 & \sum_i F_i K_i + \sum_{h,i} V_i q_{i,h}
 + \sum_h \lambda_h\!\left(L_h - \sum_i q_{i,h}\right) \\
 & + \sum_{i,h} \phi_{i,h}\!\left(K_i CF_{i,h} - q_{i,h}\right)
 + \mu\!\left(\sum_i \mathrm{ELCC}_i K_i - \mathrm{PRM}\cdot\mathrm{PeakLoad}\right).
\end{aligned}
\end{equation}
The first-order conditions for $q_{i,h}$ and $K_i$ yield, for technologies
that are marginal in the optimum,
\begin{align}
V_i - \lambda_h - \phi_{i,h} & = 0, \qquad &&\text{(marginal energy cost)} \\
F_i + \sum_h \phi_{i,h} CF_{i,h} + \mu\,\mathrm{ELCC}_i & = 0,
   &&\text{(marginal capacity cost).}
\end{align}
The solution of this static problem defines:
(i) an \emph{optimal capacity vector} $K_i^{\star}$, and
(ii) an \emph{effective annualised cost} of capacity for each technology,
\begin{equation}
C_i^{\mathrm{eff}}
\;\equiv\;
F_i + \sum_h \phi_{i,h} CF_{i,h} + \mu\,\mathrm{ELCC}_i,
\end{equation}
which incorporates both system-integration and reliability burdens
via the shadow prices $\phi_{i,h}$ and $\mu$.
In the reduced-form annual model below, we treat
$\{K_i^{\star}, C_i^{\mathrm{eff}}\}$ as exogenous outputs of the
fully optimising benchmark.
\medskip

\subsection{ One-period capital updating rule.}

In the annual reduced-form model we no longer resolve hourly dispatch.
Instead, we track installed capacities $K_i^t$ at the start of year $t$,
exogenous retirements $R_i^t$ (normal end-of-life), and new investment
$I_i^t$. Capital evolves as
\begin{equation}
K_i^{t+1} = K_i^t - R_i^t + I_i^t,
\end{equation}
with no early retirements: only $R_i^t$ leaves the system.

Let net available capacity after retirements be
$\tilde{K}_i^t \equiv K_i^t - R_i^t$, and suppose demand plus the
planning reserve margin in year $t\!+\!1$ requires total firm capacity
$K_{\mathrm{req}}^{t+1}$. The additional firm capacity needed is
\begin{equation}
\Delta K_{\mathrm{req}}^t
\;=\;
\max\Bigl\{
0,\;
K_{\mathrm{req}}^{t+1} - \sum_i \mathrm{ELCC}_i \tilde{K}_i^t
\Bigr\}.
\end{equation}
We choose non-negative investments $I_i^t$ to meet exactly this requirement:
\begin{equation}
\sum_i \mathrm{ELCC}_i I_i^t = \Delta K_{\mathrm{req}}^t,
\qquad I_i^t \ge 0.
\end{equation}

To avoid over-investing in technologies that are already above their
static least-cost level, we define the remaining ``gap'' to the benchmark
capacity:
\begin{equation}
\mathrm{gap}_i^t \equiv \max\{0, K_i^{\star} - \tilde{K}_i^t\}.
\end{equation}
New investment is then allocated in two steps:

\emph{(i) Move toward the least-cost mix.}
If the required addition is smaller than the total remaining gap,
\begin{equation}
\Delta K_{\mathrm{req}}^t \le \sum_j \mathrm{ELCC}_j\,\mathrm{gap}_j^t,
\end{equation}
we distribute investment proportionally to these gaps (in firm-capacity terms):
\begin{equation}
I_i^t
\;=\;
\Delta K_{\mathrm{req}}^t
\;\frac{\mathrm{ELCC}_i\,\mathrm{gap}_i^t}
     {\sum_j \mathrm{ELCC}_j\,\mathrm{gap}_j^t}.
\end{equation}
In this case all new capacity moves the system strictly closer to the
static optimum $K_i^{\star}$, and no technology overshoots its benchmark.

\emph{(ii) Fill any remaining requirement at least cost.}
If instead
\begin{equation}
\Delta K_{\mathrm{req}}^t > \sum_j \mathrm{ELCC}_j\,\mathrm{gap}_j^t,
\end{equation}
we first close all gaps ($I_i^t \ge \mathrm{gap}_i^t$ where needed), and
allocate the remaining required firm capacity
$\Delta K_{\mathrm{req}}^t - \sum_j \mathrm{ELCC}_j\,\mathrm{gap}_j^t$
by minimising the effective cost of new build:
\begin{equation}
\min_{\{I_i^t\}} \;\sum_i C_i^{\mathrm{eff}} I_i^t
\quad \text{s.t.} \quad
\sum_i \mathrm{ELCC}_i I_i^t
=
\Delta K_{\mathrm{req}}^t - \sum_j \mathrm{ELCC}_j\,\mathrm{gap}_j^t,
\;\; I_i^t \ge 0.
\end{equation}
The solution is a simple merit-order rule: any remaining capacity beyond
$K_i^{\star}$ is built in the technologies with the lowest $C_i^{\mathrm{eff}}$.

Under this updating rule, the annual model always:
(i) installs just enough new capacity to meet demand (plus reserves),
(ii) moves monotonically toward the least-cost capacity vector
$K_i^{\star}$ whenever possible, and
(iii) uses the shadow-price-based effective costs $C_i^{\mathrm{eff}}$
to govern which technologies expand when the system is already close to
the static optimum.

\subsection{Early retirement as an extension.}

We now extend the annual model to allow for \emph{early retirement} of
existing capacity, in addition to exogenous end-of-life retirements.
Let $E_i^t \ge 0$ denote early retirements of technology $i$ in year $t$.
Capital then evolves as
\begin{equation}
K_i^{t+1} = K_i^t - R_i^t - E_i^t + I_i^t,
\end{equation}
where $R_i^t$ are normal retirements and $I_i^t$ are new investments.
We continue to rule out negative investment or re–entry:
$E_i^t \ge 0$, $I_i^t \ge 0$.

For each technology $i$, let $\bar{C}_i^{\mathrm{keep}}$ denote the
annualised cost of keeping one unit of existing capacity online for one
more year (including fixed O\&M and expected operating costs valued at
the optimiser’s prices). Let $C_i^{\mathrm{new}}$ denote the
effective annualised cost of building one unit of new capacity,
as defined in the static benchmark above:
$C_i^{\mathrm{new}} \equiv C_i^{\mathrm{eff}}$.

Given demand and a planning reserve margin in year $t\!+\!1$, total firm
capacity must again satisfy
\begin{equation}
\sum_i \mathrm{ELCC}_i K_i^{t+1} \;\ge\; K_{\mathrm{req}}^{t+1}.
\end{equation}
Substituting the capital-update equation and rearranging yields a
requirement on the net effect of early retirements and new investment:
\begin{equation}
\sum_i \mathrm{ELCC}_i\bigl(-E_i^t + I_i^t\bigr)
\;\ge\;
\Delta K_{\mathrm{req}}^t,
\end{equation}
where
\begin{equation}
\Delta K_{\mathrm{req}}^t
\;\equiv\;
K_{\mathrm{req}}^{t+1}
-
\sum_i \mathrm{ELCC}_i (K_i^t - R_i^t)
\end{equation}
is the net additional firm capacity required relative to
the post–normal-retirement stock.

Taking the static capacities $K_i^t$ as given, the annual planner chooses
early retirements and new investments to minimise the cost of meeting
this firm-capacity requirement:
\begin{equation}
\min_{\{E_i^t, I_i^t\}}
\;\sum_i \bar{C}_i^{\mathrm{keep}} (K_i^t - R_i^t - E_i^t)
   + \sum_i C_i^{\mathrm{new}} I_i^t
\quad
\text{s.t.}\quad
\sum_i \mathrm{ELCC}_i(-E_i^t + I_i^t)
\ge \Delta K_{\mathrm{req}}^t,
\;
E_i^t, I_i^t \ge 0.
\end{equation}
Dropping constant terms that do not depend on $E_i^t$ or $I_i^t$, this is
equivalent to
\begin{equation}
\min_{\{E_i^t, I_i^t\}}
\;\sum_i\Bigl(
-\,\bar{C}_i^{\mathrm{keep}} E_i^t
+ C_i^{\mathrm{new}} I_i^t
\Bigr)
\quad
\text{s.t.}\quad
\sum_i \mathrm{ELCC}_i(-E_i^t + I_i^t)
\ge \Delta K_{\mathrm{req}}^t.
\end{equation}
Let $\lambda^t \ge 0$ be the shadow price of the firm-capacity
requirement. The Lagrangian is
\begin{equation}
\mathcal{L}^t
=
\sum_i\Bigl(
-\,\bar{C}_i^{\mathrm{keep}} E_i^t
+ C_i^{\mathrm{new}} I_i^t
\Bigr)
+
\lambda^t
\Bigl(
\Delta K_{\mathrm{req}}^t
-
\sum_i \mathrm{ELCC}_i(-E_i^t + I_i^t)
\Bigr).
\end{equation}
For technologies at an interior solution ($E_i^t>0$ or $I_i^t>0$), the
first-order conditions are
\begin{align}
\frac{\partial \mathcal{L}^t}{\partial E_i^t}
&=
-\,\bar{C}_i^{\mathrm{keep}}
+ \lambda^t \mathrm{ELCC}_i
= 0
\quad\Rightarrow\quad
\lambda^t
=
\frac{\bar{C}_i^{\mathrm{keep}}}{\mathrm{ELCC}_i},
\\[0.5em]
\frac{\partial \mathcal{L}^t}{\partial I_i^t}
&=
C_i^{\mathrm{new}}
- \lambda^t \mathrm{ELCC}_i
= 0
\quad\Rightarrow\quad
\lambda^t
=
\frac{C_i^{\mathrm{new}}}{\mathrm{ELCC}_i}.
\end{align}
Thus, in optimum, the planner compares the \emph{cost per unit of firm
capacity} of keeping versus replacing:
\begin{equation}
\tilde{C}_i^{\mathrm{keep}}
\;\equiv\;
\frac{\bar{C}_i^{\mathrm{keep}}}{\mathrm{ELCC}_i},
\qquad
\tilde{C}_i^{\mathrm{new}}
\;\equiv\;
\frac{C_i^{\mathrm{new}}}{\mathrm{ELCC}_i}.
\end{equation}
The optimal early-retirement rule can be stated simply as:
\begin{equation}
\text{Retire early any capacity for which }
\tilde{C}_i^{\mathrm{keep}}
>\min_j \tilde{C}_j^{\mathrm{new}},
\end{equation}
i.e.\ it is cheaper (per unit of firm capacity) to replace that plant
with the lowest-cost available new technology than to keep it online.

Once optimal early retirements $E_i^t$ have been determined using this
rule, the remaining capacity
\(
\tilde{K}_i^t = K_i^t - R_i^t - E_i^t
\)
defines a new net requirement
\begin{equation}
\Delta K_{\mathrm{req}}^{t,\mathrm{post}}
=
\max\Bigl\{
0,\;
K_{\mathrm{req}}^{t+1}
-
\sum_i \mathrm{ELCC}_i \tilde{K}_i^t
\Bigr\},
\end{equation}
and \emph{the new investment $I_i^t$ is then allocated exactly as in the
previous subsection}, using the same least-cost rule based on
$C_i^{\mathrm{new}} = C_i^{\mathrm{eff}}$. Early retirement therefore
operates as a preliminary screening step, potentially increasing the
required volume of new investment but leaving the merit order over
technologies unchanged.

\newpage


%%%%%%%%%%%%%%%%%%%%%%%%%%%%%%%%%%%%%%%%%%%%%%%%%%%%
% SYSTEM COSTS, OPTIMAL CAPACITIES, AND AN ADJUSTED VELOCITY RULE
%%%%%%%%%%%%%%%%%%%%%%%%%%%%%%%%%%%%%%%%%%%%%%%%%%%%

\section{System Costs, Optimal Capacities, and an Adjusted Velocity Rule}

In CPAT at present new investment is allocated according to the modified (multilogit) least cost approach. But this takes no account if existing capacity is optimal or not.

In this note we sketch how information on marginal system costs as a function of capacity can be
used to adjust the velocity-based investment rule. The objective is deliberately simple: we first
characterise an ``optimal'' capacity mix implied by system costs alone, and then show how the
velocity rule can be modified so that investment shares respond not only to levelised costs but
also to the gap between actual and optimal capacities. later then we formalate the approach without access to the 'ideal' capacities, only on the systems costs.

\subsection*{Optimal capacities from marginal system costs}

Let $K_i$ denote the installed capacity of technology $i$ (in MW, or normalised units), and let
$MSC_i(K_i)$ denote its marginal \emph{system} cost as a function of capacity. This marginal cost
may be thought of as the derivative of a system cost function
\[
    SC_i(K_i) = \int_0^{K_i} MSC_i(x)\,dx,
\]
which aggregates, for example, integration costs, balancing costs, network costs and any
capacity-related system benefits.

In a purely system-cost-driven optimum (holding demand fixed and abstracting from levelised
generation costs), the capacity mix $\{K_i^*\}$ satisfies the usual first-order condition:
\begin{equation}
    MSC_i(K_i^*) = \mu^{\text{sys}}
    \quad \text{for all technologies $i$ with } K_i^* > 0,
    \label{eq:MSC_opt_condition}
\end{equation}
where $\mu^{\text{sys}}$ is the common marginal system cost at the optimum. Equation
\eqref{eq:MSC_opt_condition} simply states that, at the system optimum, the marginal system
cost of expanding any active technology is equalised.

We do not need to solve explicitly for $K_i^*$ in CPAT. It is sufficient to assume that (from
engineering studies or a separate calculation) a benchmark capacity vector $K_i^*$ consistent
with \eqref{eq:MSC_opt_condition} is available, together with approximate slopes
\[
    m_i = \frac{d\,MSC_i}{dK_i}\Big|_{K_i^*},
\]
which describe how marginal system cost increases (or decreases) when capacity deviates from its
benchmark.

\subsection*{A simple investment problem with capacity-gap adjustments}

Let $s_i$ denote the share of new investment going to technology $i$ in the current year, with
$\sum_i s_i = 1$. We now define a deliberately simplified myopic investment problem:
\begin{equation}
    \min_{\{s_i\}} \sum_i
    \left[
        C_i s_i
      + q_i (K_i - K_i^*) s_i
      + \frac{\lambda}{2} s_i^2
    \right]
    \quad \text{s.t.} \quad \sum_i s_i = 1,\; s_i \ge 0,
    \label{eq:capacity_gap_objective}
\end{equation}
where $C_i$ is the base (private) LCOE, $\lambda>0$ is the familiar velocity penalty parameter,
and $q_i$ is a parameter that links the capacity gap $(K_i - K_i^*)$ to an incremental cost or
benefit in the investment decision.

Intuitively, if $K_i > K_i^*$ and $q_i > 0$, then the term $q_i (K_i - K_i^*) s_i$ raises the
effective cost of investing further in technology $i$, reflecting the fact that its capacity is
already above the system optimum. Conversely, if $K_i < K_i^*$, the same term reduces the
effective cost, encouraging investment to close the gap.

\subsection*{Lagrangian and adjusted velocity rule}

The Lagrangian for \eqref{eq:capacity_gap_objective} is
\begin{equation}
    \mathcal{L}(s,\mu)
    = \sum_i \left[
        C_i s_i + q_i (K_i - K_i^*) s_i + \frac{\lambda}{2} s_i^2
      \right]
      - \mu \left( \sum_i s_i - 1 \right),
\end{equation}
where $\mu$ is the multiplier on the share constraint. The first-order condition for an interior
solution is
\begin{equation}
    C_i + q_i (K_i - K_i^*)
    + \lambda s_i
    - \mu = 0.
\end{equation}
Define an “effective” cost
\begin{equation}
    \tilde{C}_i
    = C_i + q_i (K_i - K_i^*),
    \label{eq:effective_cost_def}
\end{equation}
which incorporates both the base LCOE and the capacity-gap adjustment. Then the
first-order condition simplifies to
\begin{equation}
    \tilde{C}_i + \lambda s_i - \mu = 0,
\end{equation}
so that
\begin{equation}
    s_i = \frac{\mu - \tilde{C}_i}{\lambda}.
\end{equation}

Summing over $i$ and using $\sum_i s_i = 1$ yields
\[
    1 = \frac{N\mu - \sum_i \tilde{C}_i}{\lambda},
\]
so that
\[
    \mu = \tilde{\bar{C}} + \frac{\lambda}{N},
    \qquad
    \tilde{\bar{C}} = \frac{1}{N}\sum_i \tilde{C}_i.
\]
Substituting back, we obtain the \emph{adjusted velocity rule}:
\begin{equation}
    s_i
    = \frac{1}{N}
      + \frac{\tilde{\bar{C}} - \tilde{C}_i}{\lambda},
    \label{eq:adjusted_velocity_rule}
\end{equation}
with $\tilde{C}_i$ defined in \eqref{eq:effective_cost_def}. The non-negativity constraint is
handled as before by clipping negative shares to zero and renormalising.

In words, the optimal investment rule remains a velocity-style linear mapping from costs to
shares, but the relevant cost is now an effective cost that reflects both the base LCOE and the
deviation of capacity from its system-optimal level. If $K_i > K_i^*$, the effective cost is higher,
and $s_i$ is smaller than in the pure velocity model. If $K_i < K_i^*$, the effective cost is lower,
and $s_i$ is larger.

\subsection*{Capacity dynamics and steady state}

To close the loop, we need a simple mapping from investment shares to capacity. A deliberately
oversimplified choice is
\begin{equation}
    K_{i,t+1}
    = K_{i,t} + \alpha (s_{i,t} - k_{i,t}),
    \label{eq:capacity_update_rule}
\end{equation}
where $k_{i,t} = K_{i,t} / \sum_j K_{j,t}$ is the capacity share of technology $i$, and $\alpha$ is a
small adjustment parameter. If total capacity is normalised or grows proportionally across
technologies, we can approximate $k_{i,t} \approx K_{i,t}$ and write
\[
    K_{i,t+1} \approx K_{i,t} + \alpha (s_{i,t} - K_{i,t}),
\]
which moves capacity toward the investment share. In steady state, $K_{i,t+1} = K_{i,t}$ implies
$K_i = s_i$, so that the steady-state capacity shares and investment shares coincide. If the
system-cost benchmark satisfies $K_i^* = s_i^*$ for some target allocation $s^*$, and the
effective costs $\tilde{C}_i$ are calibrated so that \eqref{eq:adjusted_velocity_rule} yields
$s^*$ when $K_i = K_i^*$, then the pair $(K^*, s^*)$ constitutes a fixed point of the combined
investment–capacity dynamics.

\subsection*{Summary of the algorithmic adjustment}

Under this deliberately simplified framework, a CPAT-style algorithm that uses system cost
information can proceed as follows:

\begin{enumerate}
    \item Given marginal system cost functions $MSC_i(K_i)$, identify a benchmark capacity mix
          $K_i^*$ (from engineering studies or a separate optimisation).
    \item At each time step, observe current capacities $K_{i,t}$ and compute the capacity gaps
          $(K_{i,t} - K_i^*)$.
    \item Construct effective costs
          \[
              \tilde{C}_{i,t}
              = C_{i,t} + q_i (K_{i,t} - K_i^*),
          \]
          where $q_i$ is calibrated from the slope of $MSC_i$ around $K_i^*$.
    \item Apply the adjusted velocity rule
          \[
              s_{i,t}
              = \frac{1}{N}
                + \frac{\tilde{\bar{C}}_t - \tilde{C}_{i,t}}{\lambda},
          \]
          with clipping and renormalisation as needed.
    \item Update capacities using a simple rule such as
          \[
              K_{i,t+1} = K_{i,t} + \alpha (s_{i,t} - K_{i,t}),
          \]
          or another suitable mapping consistent with CPAT’s aggregate capacity dynamics.
\end{enumerate}

Although highly stylised, this scheme shows how system-cost-derived optimal capacities can be
integrated into a velocity-style investment rule using only linear-quadratic adjustments that
preserve analytical transparency and spreadsheet implementability.

%%%%%%%%%%%%%%%%%%%%%%%%%%%%%%%%%%%%%%%%%%%%%%%%%%%%%%%%%%%%%%%%%%%%%%%%%%%%%
% ADDITIONAL SECTION: SYSTEM-COST–DRIVEN UPDATE RULE WITHOUT OPTIMAL SHARES
%%%%%%%%%%%%%%%%%%%%%%%%%%%%%%%%%%%%%%%%%%%%%%%%%%%%%%%%%%%%%%%%%%%%%%%%%%%%%

\section{System-Cost–Driven Investment Rule Without Access to Optimal Capacities}

The preceding analysis assumed that the benchmark optimal capacities $\{K_i^*\}$ were known,
allowing the construction of effective costs based on capacity gaps. In many practical settings, 
including CPAT applications, such optimal capacities may not be available. Instead, we may only
observe (or estimate) the \emph{marginal system cost} profiles of technologies as a function of
their existing capacities. In this section we derive a simple, fully self-contained adjustment to the
velocity rule that operates directly on system cost signals without requiring explicit knowledge of
$(K_i^*)$.

\subsection*{Stylised system cost structure}

Suppose system costs for each technology can be represented as follows:
\begin{equation}
    MSC_i(K_i)
    =
    \begin{cases}
        \gamma_C > 0,                      & \text{if } i = \text{coal}, \\[6pt]
        -\gamma_G < 0,                     & \text{if } i = \text{gas},  \\[6pt]
        \rho_i K_i^2, \;\rho_i > 0,        & \text{if } i \in \{\text{solar, wind, hydro}\}.
    \end{cases}
    \label{eq:stylised_MSCs}
\end{equation}

This functional form captures three common system effects:

- Coal imposes a system \emph{cost} (e.g. inflexibility, pollutant externalities) independent of its
  level, represented by the constant $\gamma_C$.

- Gas provides a system \emph{benefit} (flexibility, ramping capability), represented by the
  negative marginal cost $-\gamma_G$.

- Variable renewables (solar, wind, hydro) impose increasing system costs as their penetration
  grows, captured by a simple quadratic marginal cost $\rho_i K_i^2$.

Crucially, (\ref{eq:stylised_MSCs}) can be evaluated using only current capacity levels, meaning
that no optimal solution $(K_i^*)$ is required.

\subsection*{Total system-adjusted marginal cost}

To embed these system costs into the investment rule, we define an adjusted marginal cost for
technology $i$:
\begin{equation}
    \widehat{C}_i(K_i)
    = C_i + \eta \, MSC_i(K_i),
    \label{eq:adjusted_cost_no_Kstar}
\end{equation}
where $C_i$ is the base LCOE and $\eta$ is a scaling parameter converting system costs (which
may be in \$/MW or \$/MW-year) into units compatible with levelised energy cost.

Thus,
\[
    \widehat{C}_{\text{coal}} = C_{\text{coal}} + \eta \gamma_C,
\]
\[
    \widehat{C}_{\text{gas}} = C_{\text{gas}} - \eta \gamma_G,
\]
\[
    \widehat{C}_{\text{res},i} = C_{\text{res},i} + \eta \rho_i K_i^2.
\]

These adjusted costs reflect the real-time system value/opportunity cost of further investment
in each technology.

\subsection*{Lagrangian formulation and first-order condition}

As before, we consider a velocity-style quadratic objective:
\begin{equation}
    \min_{\{s_i\}}
    \sum_i \left[
        \widehat{C}_i(K_i) s_i
      + \frac{\lambda}{2} s_i^2
    \right]
    \quad \text{s.t.} \quad \sum_i s_i = 1,\; s_i \ge 0.
    \label{eq:system_cost_velocity_objective}
\end{equation}

The Lagrangian is
\begin{equation}
    \mathcal{L}(s,\mu)
    = \sum_i \left[
        \widehat{C}_i(K_i) s_i
      + \frac{\lambda}{2} s_i^2
    \right]
    - \mu \left( \sum_i s_i - 1 \right).
\end{equation}

The first-order condition for $s_i$ is unchanged in structure:
\begin{equation}
    \widehat{C}_i(K_i) + \lambda s_i - \mu = 0,
\end{equation}
leading to
\begin{equation}
    s_i = \frac{\mu - \widehat{C}_i(K_i)}{\lambda}.
\end{equation}

Summing over all $i$ yields the same expression for the multiplier:
\[
    \mu = \widehat{\bar{C}} + \frac{\lambda}{N},
    \qquad
    \widehat{\bar{C}} = \frac{1}{N} \sum_i \widehat{C}_i(K_i).
\]

\subsection*{Updated investment rule}

The resulting update rule is:
\begin{equation}
    s_i
    = \frac{1}{N}
      + \frac{\widehat{\bar{C}} - \widehat{C}_i(K_i)}{\lambda},
    \label{eq:system_cost_adjusted_velocity_rule}
\end{equation}
with non-negativity enforced by clipping and renormalisation.

Since $\widehat{C}_i(K_i)$ incorporates the system effects directly, the rule automatically reduces
investment in coal due to the penalty $\gamma_C$, increases investment in gas due to the system
benefit $-\gamma_G$, and moderates investment in renewables when their existing capacity is
high due to the quadratic system burden $\rho_i K_i^2$.

\subsection*{Capacity update and implicit convergence}

If capacities evolve according to the simple update rule
\begin{equation}
    K_{i,t+1} = K_{i,t} + \alpha (s_{i,t} - K_{i,t}),
    \label{eq:no_Kstar_capacity_update}
\end{equation}
then (\ref{eq:system_cost_adjusted_velocity_rule}) creates a self-contained dynamic system:
capacities determine system costs, which determine adjusted costs, which determine shares,
which update capacities.

A steady state satisfies
\[
    K_i = s_i,
\]
and
\[
    \widehat{C}_i(K_i) = \widehat{\bar{C}}.
\]
Thus the system converges (if stable) to a configuration in which the marginal system-adjusted
costs of all active technologies are equalised—even though no explicit optimal capacities
$\{K_i^*\}$ were ever required.

\subsection*{Interpretation}

This extension demonstrates that the velocity model can incorporate system cost information
even when optimal capacities are unavailable. The system-cost functional forms in
(\ref{eq:stylised_MSCs}) play an analogous role to the gap term $(K_i - K_i^*)$ in the previous
section: they modify each technology’s effective cost based on its current penetration in the
system. The resulting rule remains analytically transparent, computationally light, and easily
implemented in CPAT.

%%%%%%%%%%%%%%%%%%%%%%%%%%%%%%%%%%%%%%%%%%%%%%%%%%%%%%%%%%%%%%%%%%%%%%%%%%%%%
% END OF ADDITIONAL SECTION
%%%%%%%%%%%%%%%%%%%%%%%%%%%%%%%%%%%%%%%%%%%%%%%%%%%%%%%%%%%%%%%%%%%%%%%%%%%%%

\section{Mathematical Derivations of Velocity and Acceleration Model Quantities}

This appendix provides derivations omitted from the main text for readability.

\subsection{Derivation of the velocity model share rule}

Consider the problem:
\[
    \min_{\{s_i\}} \sum_i \left( C_i s_i + \frac{\lambda}{2} s_i^2 \right)
    \quad \text{s.t.} \quad \sum_i s_i = 1.
\]
Form the Lagrangian:
\[
    \mathcal{L} = \sum_i \left( C_i s_i + \frac{\lambda}{2} s_i^2 \right)
    - \mu \left( \sum_i s_i - 1 \right).
\]
Taking first-order conditions yields:
\[
    C_i + \lambda s_i - \mu = 0,
\]
so that
\[
    s_i = \frac{\mu - C_i}{\lambda}.
\]
Summing over $i$ produces:
\[
    \sum_i s_i = \frac{N\mu - \sum_i C_i}{\lambda} = 1,
\]
which implies:
\[
    \mu = \bar{C} + \frac{\lambda}{N}.
\]
Substituting back yields:
\[
    s_i = \frac{1}{N} + \frac{\bar{C} - C_i}{\lambda},
\]
as given in the main text.

\subsection{Derivation of the acceleration model update rule}

Consider the one-period problem:
\[
    \min_{\{s_{i,t}\}}
    \sum_i \left( C_{i,t} s_{i,t} + \frac{\lambda}{2} (s_{i,t}-s_{i,t-1})^2 \right),
\]
subject to $\sum_i s_{i,t} = 1$. Form the Lagrangian:
\[
    \mathcal{L}_t = \sum_i \left( C_{i,t} s_{i,t}
    + \frac{\lambda}{2} (s_{i,t} - s_{i,t-1})^2 \right)
    - \mu_t \left( \sum_i s_{i,t} - 1 \right).
\]
Differentiating yields:
\[
    C_{i,t} + \lambda (s_{i,t} - s_{i,t-1}) - \mu_t = 0.
\]
Thus,
\[
    s_{i,t} = s_{i,t-1} + \frac{\mu_t - C_{i,t}}{\lambda}.
\]
Summing over $i$ yields:
\[
    \mu_t = \frac{1}{N} \sum_i C_{i,t}.
\]


\subsection{Alternative Penalty Structures}

The penalty framework can be generalised in several directions. For instance, replacing the
quadratic penalty with a power penalty:
\[
    \gamma |s_i|^p,
    \quad p > 1,
\]
allows greater control over the curvature of the diversification effect. Alternatively, an entropy
penalty:
\[
    - \eta \sum_i s_i \ln s_i
\]
yields an allocation similar to the logit model but permits zero shares when $\eta$ is small. These
extensions may be useful where the aim is to mimic optimisation models with explicit
diversification constraints.

%%%%%%%%%%%%%%%%%%%%%%%%%%%%%%%%%%%%%%%%%%%%%%%%%%%%
% APPENDIX C: ALGORITHMIC IMPLEMENTATION
%%%%%%%%%%%%%%%%%%%%%%%%%%%%%%%%%%%%%%%%%%%%%%%%%%%%


\newpage
\newpage
\section{Algorithmic Implementation}

This appendix describes a practical algorithm for implementing the combined velocity–acceleration allocation rule in a spreadsheet environment such as Excel. The simple models and the combined model lend themselves to simple algorithmic implementation suitable for spreadsheets.

\subsection{Background: Simple Velocity model}

\begin{enumerate}
    \item Compute the cost average $\bar{C}$.
    \item Evaluate $s_i = \frac{1}{N} + \frac{\bar{C} - C_i}{\lambda}$.
    \item Set negative values to zero.
    \item Renormalise to ensure $\sum_i s_i = 1$.
\end{enumerate}

\subsection{Background: Simple Acceleration model}

\begin{enumerate}
    \item Compute $\bar{C}_t$.
    \item Update shares using $s_{i,t} = s_{i,t-1} + (\bar{C}_t - C_{i,t})/\lambda$.
    \item Clip negative shares and renormalise.
\end{enumerate}


\subsection{Implementation In Excel}

The starting
point is the closed-form update equation derived in Section~6,
\begin{equation}
    s_{i,t}
    = \frac{1}{N}
      + \frac{
            \bar{C}_t^{\text{adj}}
          - C_{i,t}^{\text{base}}
          - C_{i,t}^{\text{sys}}(s_{t-1};D)
          + \lambda_a s_{i,t-1}
        }{\lambda_v + \lambda_a},
    \label{eq:A3_combined_update_final}
\end{equation}
where the adjusted average cost is
\begin{equation}
    \bar{C}_t^{\text{adj}}
    = \bar{C}_t^{\text{base}} + \bar{C}_t^{\text{sys}}(s_{t-1};D) - \frac{\lambda_a}{N}.
    \label{eq:A3_Cbar_adj}
\end{equation}
Here $\lambda_v$ is the velocity penalty, $\lambda_a$ is the acceleration penalty (typically set
smaller), $N$ is the number of technologies, and $D$ is a country-level population density
measure.

The base cost and system adjustments are defined as
\begin{equation}
    C_{i,t}^{\text{tot}}(s_t,s_{t-1};D)
    = C_{i,t}^{\text{base}}
      + C_{i,t}^{\text{sys}}(s_{t-1};D)
      + \lambda_v s_{i,t}
      + \lambda_a (s_{i,t} - s_{i,t-1}),
\end{equation}
with system costs
\begin{equation}
    C_{i,t}^{\text{sys}}(s_{t-1};D)
    = - \beta_G G_i s_{\text{gas},t-1}
      + \kappa_R R_i S_{\text{RES},t-1}
      + \frac{\phi_i}{D} s_{i,t-1},
    \label{eq:A3_Csys}
\end{equation}
and
\[
    S_{\text{RES},t-1} = \sum_j R_j s_{j,t-1}.
\]

In a spreadsheet, the implementation proceeds year by year. The following steps describe the
algorithm for period $t$, given known values of costs, parameters and previous-year shares
$s_{i,t-1}$.

\paragraph{Step 1: Set up inputs.}

For each technology $i$ and year $t$, the spreadsheet should contain:

\begin{itemize}
    \item the base LCOE, $C_{i,t}^{\text{base}}$;
    \item technology indicators $G_i$ (gas), $R_i$ (variable renewables);
    \item previous-year shares $s_{i,t-1}$;
    \item common parameters $\lambda_v$, $\lambda_a$, $\beta_G$, $\kappa_R$, and $\phi_i$;
    \item population density $D$.
\end{itemize}

These can be arranged in a standard technology-by-year matrix, with separate columns for each
parameter.

\paragraph{Step 2: Compute aggregate renewable share and system adjustments.}

First compute the aggregate share of variable renewables in year $t-1$:
\[
    S_{\text{RES},t-1} = \sum_j R_j s_{j,t-1}.
\]
In Excel this can be implemented as a SUMPRODUCT over the row or column containing $R_j$ and
$s_{j,t-1}$.

Then compute the system adjustment term for each technology using
(\ref{eq:A3_Csys}):
\[
    C_{i,t}^{\text{sys}}(s_{t-1};D)
    = - \beta_G G_i s_{\text{gas},t-1}
      + \kappa_R R_i S_{\text{RES},t-1}
      + \frac{\phi_i}{D} s_{i,t-1}.
\]
This requires (i) reading $s_{\text{gas},t-1}$ from the previous-year gas share, (ii) multiplying
the renewables indicator $R_i$ by $S_{\text{RES},t-1}$, and (iii) scaling the previous-year share
by $\phi_i / D$.

\paragraph{Step 3: Compute average base and system costs.}

Compute the simple averages:
\[
    \bar{C}_t^{\text{base}} = \frac{1}{N} \sum_i C_{i,t}^{\text{base}}, 
    \qquad
    \bar{C}_t^{\text{sys}}(s_{t-1};D) = \frac{1}{N} \sum_i C_{i,t}^{\text{sys}}(s_{t-1};D),
\]
using an AVERAGE or SUM/N formula across technologies.

Then compute the adjusted average cost as in (\ref{eq:A3_Cbar_adj}):
\[
    \bar{C}_t^{\text{adj}}
    = \bar{C}_t^{\text{base}} + \bar{C}_t^{\text{sys}}(s_{t-1};D) - \frac{\lambda_a}{N}.
\]

\paragraph{Step 4: Compute preliminary new shares (before clipping).}

For each technology, compute the unconstrained share $s_{i,t}^{\text{raw}}$ using
(\ref{eq:A3_combined_update_final}):
\[
    s_{i,t}^{\text{raw}}
    = \frac{1}{N}
      + \frac{
            \bar{C}_t^{\text{adj}}
          - C_{i,t}^{\text{base}}
          - C_{i,t}^{\text{sys}}(s_{t-1};D)
          + \lambda_a s_{i,t-1}
        }{\lambda_v + \lambda_a}.
\]
In Excel, this is a single cell formula referencing the relevant row values for $C_{i,t}^{\text{base}}$,
$C_{i,t}^{\text{sys}}$, $s_{i,t-1}$ and the scalar parameters $\lambda_v$, $\lambda_a$ and $N$.

\paragraph{Step 5: Apply non-negativity and renormalise.}

To enforce $s_{i,t} \ge 0$, define:
\[
    \tilde{s}_{i,t} = \max\{0, s_{i,t}^{\text{raw}}\}.
\]
In Excel this is implemented with a MAX formula, for example
\texttt{=MAX(0, s\_raw)}.

Then renormalise the shares so they sum to one:
\[
    s_{i,t} = \frac{\tilde{s}_{i,t}}{\sum_j \tilde{s}_{j,t}}.
\]
This requires computing the sum of $\tilde{s}_{j,t}$ across all technologies and dividing each
technology’s clipped share by this sum.

\paragraph{Step 6: Iterate over time.}

Once the sheet is set up for a single year, the same formulas can be copied across columns to
simulate multiple years. For each $t+1$, the previous-year shares $s_{i,t}$ become the inputs
$s_{i,t}$ in Step~1. The system automatically generates a sequence of investment shares
consistent with the combined velocity–acceleration model and the chosen parameters.

\paragraph{Summary.}

The combined model can thus be implemented with standard spreadsheet functions (SUM,
SUMPRODUCT, AVERAGE, MAX and simple arithmetic), ensuring compatibility with CPAT’s
existing Excel-based architecture. The key policy-relevant inputs are the base LCOEs,
system-cost parameters $(\beta_G,\kappa_R,\phi_i)$, the penalty parameters $(\lambda_v,
\lambda_a)$ and the population density $D$. All derived quantities, including the total effective
levelised cost and the evolution of shares, follow mechanically from the above steps.

%%%%%%%%%%%%%%%%%%%%%%%%%%%%%%%%%%%%%%%%%%%%%%%%%%%%
% END 
%%%%%%%%%%%%%%%%%%%%%%%%%%%%%%%%%%%%%%%%%%%%%%%%%%%%
%%%%%%%%%%%%%%%%%%%%%%%%%%%%%%%%%%%%%%%%%%%%%%%%%%%%%%%%%%%%%%%%%%%%%%%%%%%%%
% APPENDIX D. FINAL SUMMARY: SYSTEM-COST–BASED INVESTMENT AND RETIREMENT
%%%%%%%%%%%%%%%%%%%%%%%%%%%%%%%%%%%%%%%%%%%%%%%%%%%%%%%%%%%%%%%%%%%%%%%%%%%%%

\newpage

%%%%%%%%%%%%%%%%%%%%%%%%%%%%%%%%%%%%%%%%%%%%%%%%%%%%%%%%%%%%%%%%%%%%%%%%%%%%%
% APPENDIX D. FINAL SUMMARY: OPTIMAL SYSTEM-COST–BASED INVESTMENT & RETIREMENT
%%%%%%%%%%%%%%%%%%%%%%%%%%%%%%%%%%%%%%%%%%%%%%%%%%%%%%%%%%%%%%%%%%%%%%%%%%%%%

\section{Optimal System-Cost–Based Investment and Retirement Rules}

This appendix restates the yearly model in a simple, stylised form and derives the \emph{optimal}
retirement and investment rules under the assumption that all decisions are driven by system
costs rather than by an explicitly computed optimal capacity mix. System costs are taken as
given functions of capacity, and the allocation of new investment responds to \emph{effective}
costs using a combined velocity–acceleration rule. Retirement responds to the same effective
costs through a simple static optimisation. The framework is annual (one step per year) and does
not rely on 8760-hour operational detail.

\subsection{System costs and effective annual costs}

Let $i$ index technologies and $t$ index years. Let $K_{i,t}$ denote installed capacity of technology
$i$ at the start of year $t$, and let $C_i$ denote its base levelised cost of electricity (LCOE),
assumed constant over time for simplicity.

We assume stylised system costs per unit of capacity given by
\begin{equation}
    MSC_i(K_{i,t})
    =
    \begin{cases}
        \gamma_C > 0,                        & \text{if } i = \text{coal}, \\[4pt]
        -\gamma_G < 0,                       & \text{if } i = \text{gas},  \\[4pt]
        \rho_i K_{i,t}^2, \;\rho_i > 0,      & \text{if } i \in \{\text{solar, wind, hydro}\},
    \end{cases}
    \label{eq:D_MSCs_opt}
\end{equation}
where $\gamma_C$ represents a constant system penalty for coal, $-\gamma_G$ a constant system
benefit for gas (flexibility), and $\rho_i K_{i,t}^2$ an increasing system burden from variable
renewables as their penetration grows.

We convert these system costs into an \emph{effective} annual cost per MWh:
\begin{equation}
    C_{i,t}^{\mathrm{eff}}
    = C_i + \eta \, MSC_i(K_{i,t}),
    \label{eq:D_Ceff_opt}
\end{equation}
where $\eta$ is a scaling parameter aligning units. The effective cost
$C_{i,t}^{\mathrm{eff}}$ is the key object that drives both retirement and investment.

We also define the average effective cost across all $N$ technologies:
\begin{equation}
    \bar{C}_t^{\mathrm{eff}}
    = \frac{1}{N} \sum_{j=1}^N C_{j,t}^{\mathrm{eff}}.
    \label{eq:D_Cbar_opt}
\end{equation}

For generation, let $CF_i$ denote the (constant) capacity factor of technology $i$. Total
generation in year $t$ is
\[
    G_t = \sum_i CF_i K_{i,t},
\]
and we treat required generation $G_t^{\mathrm{req}}$ as exogenous. The system must choose
retirements and investments so that generation after retirement and investment satisfies
\[
    \sum_i CF_i K_{i,t+1} = G_t^{\mathrm{req}}.
\]


\subsection{Simplified System-Cost–Driven Retirement and Investment}

This appendix provides a concise, stylised yearly model in which investment and retirement
respond to system-adjusted effective costs. All quantities are annual; no hourly structure is used.

Let $K_{i,t}$ denote installed capacity of technology $i$ at the start of year $t$, and $C_i$ its base
LCOE. System costs are given exogenously as:
\[
MSC_i(K_{i,t}) =
\begin{cases}
\gamma_C > 0, & \text{if } i = \text{coal},\\[3pt]
\rho_i K_{i,t}^2, & \text{if } i \in \mathcal{R} \; (\text{solar, wind, hydro}).
\end{cases}
\]

Effective marginal cost (per MWh) is
\[
C_{i,t}^{\mathrm{eff}} = C_i + \eta\, MSC_i(K_{i,t}).
\]

Let $C_{t}^{\mathrm{RES,min}} = \min_{j \in \mathcal{R}} C_{j,t}^{\mathrm{eff}}$ denote the cheapest renewable
effective cost at time $t$. Coal has a (given) variable marginal cost $C_{C,t}^{\mathrm{var}}$.


\subsection{Optimal investment rule with velocity and acceleration}

After retirement, capacity is $K_{i,t}^{\mathrm{pre\text{-}inv}}$. Let the required generation in year
$t+1$ be $G_{t+1}^{\mathrm{req}}$. The total capacity addition needed to meet the future requirement
is determined externally (for example, from demand projections and capacity factors). Denote
the total new capacity by $I_t$ and its allocation across technologies by $I_{i,t}$, so that
\[
    \sum_i CF_i I_{i,t} = G_{t+1}^{\mathrm{req}} - \sum_i CF_i K_{i,t}^{\mathrm{pre\text{-}inv}}.
\]
Define $s_{i,t} = I_{i,t} / I_t$ as the \emph{share} of new capacity investment going to technology
$i$, so that $\sum_i s_{i,t} = 1$.

We now derive the optimal investment shares by solving a quadratic problem that combines
effective costs with velocity and acceleration penalties. The planner chooses $\{s_{i,t}\}$ to
minimise:
\begin{equation}
    \min_{\{s_{i,t}\}}
    \sum_i \left[
        C_{i,t}^{\mathrm{eff}} s_{i,t}
      + \frac{\lambda_v}{2} s_{i,t}^2
      + \frac{\lambda_a}{2} (s_{i,t} - s_{i,t-1})^2
    \right]
    \quad \text{s.t.} \quad \sum_i s_{i,t} = 1,
    \label{eq:D_inv_opt_problem}
\end{equation}
where $\lambda_v > 0$ is the velocity penalty parameter and $\lambda_a \ge 0$ is the acceleration
penalty parameter. This problem is fully analogous to the combined velocity–acceleration
formulation presented earlier, but now the costs entering the objective are the effective costs
$C_{i,t}^{\mathrm{eff}}$ derived from system costs.

Form the Lagrangian:
\begin{equation}
    \mathcal{L}_{\mathrm{inv}}
    = \sum_i \left[
        C_{i,t}^{\mathrm{eff}} s_{i,t}
      + \frac{\lambda_v}{2} s_{i,t}^2
      + \frac{\lambda_a}{2} (s_{i,t} - s_{i,t-1})^2
    \right]
    - \mu_t \left( \sum_i s_{i,t} - 1 \right),
\end{equation}
where $\mu_t$ enforces the share constraint. The first-order condition with respect to $s_{i,t}$
is
\begin{equation}
    C_{i,t}^{\mathrm{eff}}
    + \lambda_v s_{i,t}
    + \lambda_a (s_{i,t} - s_{i,t-1})
    - \mu_t = 0.
    \label{eq:D_inv_FOC}
\end{equation}
Collecting terms in $s_{i,t}$ gives:
\begin{equation}
    (\lambda_v + \lambda_a) s_{i,t}
    = \mu_t
      - C_{i,t}^{\mathrm{eff}}
      + \lambda_a s_{i,t-1}.
    \label{eq:D_inv_linear}
\end{equation}
Thus,
\begin{equation}
    s_{i,t}
    = \frac{1}{\lambda_v + \lambda_a}
      \left[
          \mu_t
        - C_{i,t}^{\mathrm{eff}}
        + \lambda_a s_{i,t-1}
      \right].
    \label{eq:D_inv_pre_mu}
\end{equation}

To determine $\mu_t$, sum \eqref{eq:D_inv_pre_mu} over $i$, using $\sum_i s_{i,t} = 1$ and
$\sum_i s_{i,t-1} = 1$:
\[
    1
    = \frac{N\mu_t
        - \sum_i C_{i,t}^{\mathrm{eff}}
        + \lambda_a \sum_i s_{i,t-1}
      }{\lambda_v + \lambda_a}
    = \frac{N\mu_t
        - \sum_i C_{i,t}^{\mathrm{eff}}
        + \lambda_a
      }{\lambda_v + \lambda_a}.
\]
Therefore,
\begin{equation}
    \mu_t
    = \bar{C}_t^{\mathrm{eff}}
      - \frac{\lambda_a}{N}
      + \frac{\lambda_v + \lambda_a}{N},
\end{equation}
where $\bar{C}_t^{\mathrm{eff}}$ is the average effective cost defined in \eqref{eq:D_Cbar_opt}.

Substituting back into \eqref{eq:D_inv_pre_mu} yields the closed-form optimal investment rule:
\begin{equation}
    s_{i,t}
    = \frac{1}{N}
      + \frac{
            \bar{C}_t^{\mathrm{eff}}
          - C_{i,t}^{\mathrm{eff}}
          + \lambda_a (s_{i,t-1} - \tfrac{1}{N})
        }{\lambda_v + \lambda_a}.
    \label{eq:D_inv_rule_final}
\end{equation}

When $\lambda_a = 0$, this collapses to a pure velocity rule:
\[
    s_{i,t}
    = \frac{1}{N}
      + \frac{\bar{C}_t^{\mathrm{eff}} - C_{i,t}^{\mathrm{eff}}}{\lambda_v},
\]
and when $\lambda_v$ is large relative to $\lambda_a$, the rule closely tracks $s_{i,t-1}$ with only
gradual adjustments driven by cost differentials.

In practice, negative $s_{i,t}$ values are clipped to zero and the remaining positive shares are
renormalised to sum to one.

\subsection{Capacity update with optimal retirement and investment}

After choosing optimal retirements and investments, capacity evolves according to
\begin{equation}
    K_{i,t+1}
    = K_{i,t}^{\mathrm{pre\text{-}inv}} + I_{i,t}
    = (K_{i,t} - R_{i,t}) + I_t s_{i,t},
    \label{eq:D_K_update_opt}
\end{equation}
where $R_{i,t}$ solves the retirement problem \eqref{eq:D_ret_opt_problem} (i.e.\ merit-order
retirement based on $C_{i,t}^{\mathrm{eff}}$), and $s_{i,t}$ solves the investment problem
\eqref{eq:D_inv_opt_problem} (velocity–acceleration optimum based on $C_{i,t}^{\mathrm{eff}}$).

The total investment $I_t$ is chosen (outside this appendix) to ensure that the generation
constraint
\[
    \sum_i CF_i K_{i,t+1} = G_{t+1}^{\mathrm{req}}
\]
is satisfied.

\subsection{Simplified algorithmic implementation}

The optimal yearly algorithm, in stylised form, proceeds as follows.

\paragraph{Given at the start of year $t$:}
capacities $K_{i,t}$, previous investment shares $s_{i,t-1}$,
base LCOEs $C_i$, capacity factors $CF_i$,
system-cost parameters $\gamma_C, \gamma_G, \rho_i$, scaling $\eta$,
velocity and acceleration parameters $\lambda_v, \lambda_a$,
and generation requirements $G_t^{\mathrm{req}}, G_{t+1}^{\mathrm{req}}$.

\medskip

\noindent
\paragraph{Step 1: Compute effective costs.}
\begin{itemize}
    \item Compute system costs:
    \[
        MSC_i(K_{i,t}) \text{ from \eqref{eq:D_MSCs_opt}},
    \]
    \[
        C_{i,t}^{\mathrm{eff}} = C_i + \eta\, MSC_i(K_{i,t}),
    \]
    \item Compute their average:
    \[
        \bar{C}_t^{\mathrm{eff}} = \frac{1}{N} \sum_j C_{j,t}^{\mathrm{eff}}.
    \]
\end{itemize}

\noindent
\paragraph{Step 2: Retirement (coal-only, cost-triggered with velocity smoothing).}
\begin{itemize}

    \item Compute the cheapest renewable effective cost:
    \[
        C_{t}^{\mathrm{RES,min}} = \min_{j \in \mathcal{R}} C_{j,t}^{\mathrm{eff}}.
    \]

    \item Check whether coal is economically dominated:
    \[
        C_{C,t}^{\mathrm{var}} > C_{t}^{\mathrm{RES,min}}.
    \]
    If this condition is false, coal retires only at its technical rate:  
    \[
        r_{C,t} = \delta_C.
    \]

    \item If the condition is true, compute the velocity-based additional retirement incentive:
    \[
        u_{C,t}^{\mathrm{vel}}
        =
        \min\!\left(
            u_C^{\max},
            \frac{C_{C,t}^{\mathrm{var}} - C_{t}^{\mathrm{RES,min}}}{\Lambda_R}
        \right), \qquad u_{C,t}^{\mathrm{vel}} \ge 0.
    \]

    \item Smooth the retirement change using an acceleration-like rule:
    \[
        u_{C,t}
        = (1 - \alpha_R) u_{C,t-1} + \alpha_R u_{C,t}^{\mathrm{vel}}.
    \]

    \item The total coal retirement fraction is:
    \[
        r_{C,t} = \delta_C + u_{C,t}, \qquad r_{C,t} \in [0,1].
    \]

    \item For all non-coal technologies $i \neq C$, retirement is purely technical:
    \[
        r_{i,t} = \delta_i.
    \]

    \item Apply retirement to obtain pre-investment capacities:
    \[
        K_{i,t}^{\mathrm{pre\text{-}inv}} = (1 - r_{i,t}) K_{i,t}.
    \]

\end{itemize}

\noindent
\paragraph{Step 3: Required net investment.}
\begin{itemize}
    \item Compute the generation from the pre-investment fleet:
    \[
        G_t^{\mathrm{pre\text{-}inv}} = \sum_i CF_i K_{i,t}^{\mathrm{pre\text{-}inv}}.
    \]
    \item Compute the required additional generation:
    \[
        \Delta G_t = G_{t+1}^{\mathrm{req}} - G_t^{\mathrm{pre\text{-}inv}}.
    \]
    \item Translate $\Delta G_t$ into total new capacity $I_t$ (e.g.\ using an average capacity
          factor across new investments or a separate simple rule).
\end{itemize}

\noindent
\paragraph{Step 4: Optimal investment shares (velocity–acceleration rule).}
\begin{itemize}
    \item Compute the closed-form optimal shares:
    \[
        s_{i,t}
        = \frac{1}{N}
          + \frac{
                \bar{C}_t^{\mathrm{eff}}
              - C_{i,t}^{\mathrm{eff}}
              + \lambda_a (s_{i,t-1} - \tfrac{1}{N})
            }{\lambda_v + \lambda_a},
    \]
    \item Clip negative $s_{i,t}$ to zero and renormalise so that $\sum_i s_{i,t} = 1$.
\end{itemize}

\noindent
\paragraph{Step 5: Capacity update.}
\begin{itemize}
    \item Allocate new capacity:
    \[
        I_{i,t} = I_t s_{i,t}.
    \]
    \item Update capacities:
    \[
        K_{i,t+1}
        = (K_{i,t} - R_{i,t}) + I_t s_{i,t}.
    \]
\end{itemize}

\noindent
\paragraph{Step 6: Repeat.}
\begin{itemize}
    \item Proceed to year $t+1$ with updated capacities $K_{i,t+1}$ and shares $s_{i,t}$.
\end{itemize}

\medskip

This completes the stylised optimal model. All decisions—retirement and investment—are driven
by effective costs that combine base LCOE and exogenous system cost functions. Planned retirement is
determined by exogenous exponential decay for non-coal and planned-retirement for coal. Early retirement exists only for coal if the variable cost is more than the total cost of renewables. Investment shares are determined by a velocity–acceleration optimisation, both of which are analytically transparent and compatible
with a yearly, spreadsheet-based CPAT implementation.
%%%%%%%%%%%%%%%%%%%%%%%%%%%%%%%%%%%%%%%%%%%%%%%%%%%%%%%%%%%%%%%%%%%%%%%%%%%%%

\end{document}